\documentclass[11pt]{article}

    \usepackage[breakable]{tcolorbox}
    \usepackage{parskip} % Stop auto-indenting (to mimic markdown behaviour)
    

    % Basic figure setup, for now with no caption control since it's done
    % automatically by Pandoc (which extracts ![](path) syntax from Markdown).
    \usepackage{graphicx}
    % Keep aspect ratio if custom image width or height is specified
    \setkeys{Gin}{keepaspectratio}
    % Maintain compatibility with old templates. Remove in nbconvert 6.0
    \let\Oldincludegraphics\includegraphics
    % Ensure that by default, figures have no caption (until we provide a
    % proper Figure object with a Caption API and a way to capture that
    % in the conversion process - todo).
    \usepackage{caption}
    \DeclareCaptionFormat{nocaption}{}
    \captionsetup{format=nocaption,aboveskip=0pt,belowskip=0pt}

    \usepackage{float}
    \floatplacement{figure}{H} % forces figures to be placed at the correct location
    \usepackage{xcolor} % Allow colors to be defined
    \usepackage{enumerate} % Needed for markdown enumerations to work
    \usepackage{geometry} % Used to adjust the document margins
    \usepackage{amsmath} % Equations
    \usepackage{amssymb} % Equations
    \usepackage{textcomp} % defines textquotesingle
    % Hack from http://tex.stackexchange.com/a/47451/13684:
    \AtBeginDocument{%
        \def\PYZsq{\textquotesingle}% Upright quotes in Pygmentized code
    }
    \usepackage{upquote} % Upright quotes for verbatim code
    \usepackage{eurosym} % defines \euro

    \usepackage{iftex}
    \ifPDFTeX
        \usepackage[T1]{fontenc}
        \IfFileExists{alphabeta.sty}{
              \usepackage{alphabeta}
          }{
              \usepackage[mathletters]{ucs}
              \usepackage[utf8x]{inputenc}
          }
    \else
        \usepackage{fontspec}
        \usepackage{unicode-math}
    \fi

    \usepackage{fancyvrb} % verbatim replacement that allows latex
    \usepackage{grffile} % extends the file name processing of package graphics
                         % to support a larger range
    \makeatletter % fix for old versions of grffile with XeLaTeX
    \@ifpackagelater{grffile}{2019/11/01}
    {
      % Do nothing on new versions
    }
    {
      \def\Gread@@xetex#1{%
        \IfFileExists{"\Gin@base".bb}%
        {\Gread@eps{\Gin@base.bb}}%
        {\Gread@@xetex@aux#1}%
      }
    }
    \makeatother
    \usepackage[Export]{adjustbox} % Used to constrain images to a maximum size
    \adjustboxset{max size={0.9\linewidth}{0.9\paperheight}}

    % The hyperref package gives us a pdf with properly built
    % internal navigation ('pdf bookmarks' for the table of contents,
    % internal cross-reference links, web links for URLs, etc.)
    \usepackage{hyperref}
    % The default LaTeX title has an obnoxious amount of whitespace. By default,
    % titling removes some of it. It also provides customization options.
    \usepackage{titling}
    \usepackage{longtable} % longtable support required by pandoc >1.10
    \usepackage{booktabs}  % table support for pandoc > 1.12.2
    \usepackage{array}     % table support for pandoc >= 2.11.3
    \usepackage{calc}      % table minipage width calculation for pandoc >= 2.11.1
    \usepackage[inline]{enumitem} % IRkernel/repr support (it uses the enumerate* environment)
    \usepackage[normalem]{ulem} % ulem is needed to support strikethroughs (\sout)
                                % normalem makes italics be italics, not underlines
    \usepackage{soul}      % strikethrough (\st) support for pandoc >= 3.0.0
    \usepackage{mathrsfs}
    

    
    % Colors for the hyperref package
    \definecolor{urlcolor}{rgb}{0,.145,.698}
    \definecolor{linkcolor}{rgb}{.71,0.21,0.01}
    \definecolor{citecolor}{rgb}{.12,.54,.11}

    % ANSI colors
    \definecolor{ansi-black}{HTML}{3E424D}
    \definecolor{ansi-black-intense}{HTML}{282C36}
    \definecolor{ansi-red}{HTML}{E75C58}
    \definecolor{ansi-red-intense}{HTML}{B22B31}
    \definecolor{ansi-green}{HTML}{00A250}
    \definecolor{ansi-green-intense}{HTML}{007427}
    \definecolor{ansi-yellow}{HTML}{DDB62B}
    \definecolor{ansi-yellow-intense}{HTML}{B27D12}
    \definecolor{ansi-blue}{HTML}{208FFB}
    \definecolor{ansi-blue-intense}{HTML}{0065CA}
    \definecolor{ansi-magenta}{HTML}{D160C4}
    \definecolor{ansi-magenta-intense}{HTML}{A03196}
    \definecolor{ansi-cyan}{HTML}{60C6C8}
    \definecolor{ansi-cyan-intense}{HTML}{258F8F}
    \definecolor{ansi-white}{HTML}{C5C1B4}
    \definecolor{ansi-white-intense}{HTML}{A1A6B2}
    \definecolor{ansi-default-inverse-fg}{HTML}{FFFFFF}
    \definecolor{ansi-default-inverse-bg}{HTML}{000000}

    % common color for the border for error outputs.
    \definecolor{outerrorbackground}{HTML}{FFDFDF}

    % commands and environments needed by pandoc snippets
    % extracted from the output of `pandoc -s`
    \providecommand{\tightlist}{%
      \setlength{\itemsep}{0pt}\setlength{\parskip}{0pt}}
    \DefineVerbatimEnvironment{Highlighting}{Verbatim}{commandchars=\\\{\}}
    % Add ',fontsize=\small' for more characters per line
    \newenvironment{Shaded}{}{}
    \newcommand{\KeywordTok}[1]{\textcolor[rgb]{0.00,0.44,0.13}{\textbf{{#1}}}}
    \newcommand{\DataTypeTok}[1]{\textcolor[rgb]{0.56,0.13,0.00}{{#1}}}
    \newcommand{\DecValTok}[1]{\textcolor[rgb]{0.25,0.63,0.44}{{#1}}}
    \newcommand{\BaseNTok}[1]{\textcolor[rgb]{0.25,0.63,0.44}{{#1}}}
    \newcommand{\FloatTok}[1]{\textcolor[rgb]{0.25,0.63,0.44}{{#1}}}
    \newcommand{\CharTok}[1]{\textcolor[rgb]{0.25,0.44,0.63}{{#1}}}
    \newcommand{\StringTok}[1]{\textcolor[rgb]{0.25,0.44,0.63}{{#1}}}
    \newcommand{\CommentTok}[1]{\textcolor[rgb]{0.38,0.63,0.69}{\textit{{#1}}}}
    \newcommand{\OtherTok}[1]{\textcolor[rgb]{0.00,0.44,0.13}{{#1}}}
    \newcommand{\AlertTok}[1]{\textcolor[rgb]{1.00,0.00,0.00}{\textbf{{#1}}}}
    \newcommand{\FunctionTok}[1]{\textcolor[rgb]{0.02,0.16,0.49}{{#1}}}
    \newcommand{\RegionMarkerTok}[1]{{#1}}
    \newcommand{\ErrorTok}[1]{\textcolor[rgb]{1.00,0.00,0.00}{\textbf{{#1}}}}
    \newcommand{\NormalTok}[1]{{#1}}

    % Additional commands for more recent versions of Pandoc
    \newcommand{\ConstantTok}[1]{\textcolor[rgb]{0.53,0.00,0.00}{{#1}}}
    \newcommand{\SpecialCharTok}[1]{\textcolor[rgb]{0.25,0.44,0.63}{{#1}}}
    \newcommand{\VerbatimStringTok}[1]{\textcolor[rgb]{0.25,0.44,0.63}{{#1}}}
    \newcommand{\SpecialStringTok}[1]{\textcolor[rgb]{0.73,0.40,0.53}{{#1}}}
    \newcommand{\ImportTok}[1]{{#1}}
    \newcommand{\DocumentationTok}[1]{\textcolor[rgb]{0.73,0.13,0.13}{\textit{{#1}}}}
    \newcommand{\AnnotationTok}[1]{\textcolor[rgb]{0.38,0.63,0.69}{\textbf{\textit{{#1}}}}}
    \newcommand{\CommentVarTok}[1]{\textcolor[rgb]{0.38,0.63,0.69}{\textbf{\textit{{#1}}}}}
    \newcommand{\VariableTok}[1]{\textcolor[rgb]{0.10,0.09,0.49}{{#1}}}
    \newcommand{\ControlFlowTok}[1]{\textcolor[rgb]{0.00,0.44,0.13}{\textbf{{#1}}}}
    \newcommand{\OperatorTok}[1]{\textcolor[rgb]{0.40,0.40,0.40}{{#1}}}
    \newcommand{\BuiltInTok}[1]{{#1}}
    \newcommand{\ExtensionTok}[1]{{#1}}
    \newcommand{\PreprocessorTok}[1]{\textcolor[rgb]{0.74,0.48,0.00}{{#1}}}
    \newcommand{\AttributeTok}[1]{\textcolor[rgb]{0.49,0.56,0.16}{{#1}}}
    \newcommand{\InformationTok}[1]{\textcolor[rgb]{0.38,0.63,0.69}{\textbf{\textit{{#1}}}}}
    \newcommand{\WarningTok}[1]{\textcolor[rgb]{0.38,0.63,0.69}{\textbf{\textit{{#1}}}}}
    \makeatletter
    \newsavebox\pandoc@box
    \newcommand*\pandocbounded[1]{%
      \sbox\pandoc@box{#1}%
      % scaling factors for width and height
      \Gscale@div\@tempa\textheight{\dimexpr\ht\pandoc@box+\dp\pandoc@box\relax}%
      \Gscale@div\@tempb\linewidth{\wd\pandoc@box}%
      % select the smaller of both
      \ifdim\@tempb\p@<\@tempa\p@
        \let\@tempa\@tempb
      \fi
      % scaling accordingly (\@tempa < 1)
      \ifdim\@tempa\p@<\p@
        \scalebox{\@tempa}{\usebox\pandoc@box}%
      % scaling not needed, use as it is
      \else
        \usebox{\pandoc@box}%
      \fi
    }
    \makeatother

    % Define a nice break command that doesn't care if a line doesn't already
    % exist.
    \def\br{\hspace*{\fill} \\* }
    % Math Jax compatibility definitions
    \def\gt{>}
    \def\lt{<}
    \let\Oldtex\TeX
    \let\Oldlatex\LaTeX
    \renewcommand{\TeX}{\textrm{\Oldtex}}
    \renewcommand{\LaTeX}{\textrm{\Oldlatex}}
    % Document parameters
    % Document title
    \title{Notebook\_BVP\_to\_end}
    
    
    
    
    
    
    
% Pygments definitions
\makeatletter
\def\PY@reset{\let\PY@it=\relax \let\PY@bf=\relax%
    \let\PY@ul=\relax \let\PY@tc=\relax%
    \let\PY@bc=\relax \let\PY@ff=\relax}
\def\PY@tok#1{\csname PY@tok@#1\endcsname}
\def\PY@toks#1+{\ifx\relax#1\empty\else%
    \PY@tok{#1}\expandafter\PY@toks\fi}
\def\PY@do#1{\PY@bc{\PY@tc{\PY@ul{%
    \PY@it{\PY@bf{\PY@ff{#1}}}}}}}
\def\PY#1#2{\PY@reset\PY@toks#1+\relax+\PY@do{#2}}

\@namedef{PY@tok@w}{\def\PY@tc##1{\textcolor[rgb]{0.73,0.73,0.73}{##1}}}
\@namedef{PY@tok@c}{\let\PY@it=\textit\def\PY@tc##1{\textcolor[rgb]{0.24,0.48,0.48}{##1}}}
\@namedef{PY@tok@cp}{\def\PY@tc##1{\textcolor[rgb]{0.61,0.40,0.00}{##1}}}
\@namedef{PY@tok@k}{\let\PY@bf=\textbf\def\PY@tc##1{\textcolor[rgb]{0.00,0.50,0.00}{##1}}}
\@namedef{PY@tok@kp}{\def\PY@tc##1{\textcolor[rgb]{0.00,0.50,0.00}{##1}}}
\@namedef{PY@tok@kt}{\def\PY@tc##1{\textcolor[rgb]{0.69,0.00,0.25}{##1}}}
\@namedef{PY@tok@o}{\def\PY@tc##1{\textcolor[rgb]{0.40,0.40,0.40}{##1}}}
\@namedef{PY@tok@ow}{\let\PY@bf=\textbf\def\PY@tc##1{\textcolor[rgb]{0.67,0.13,1.00}{##1}}}
\@namedef{PY@tok@nb}{\def\PY@tc##1{\textcolor[rgb]{0.00,0.50,0.00}{##1}}}
\@namedef{PY@tok@nf}{\def\PY@tc##1{\textcolor[rgb]{0.00,0.00,1.00}{##1}}}
\@namedef{PY@tok@nc}{\let\PY@bf=\textbf\def\PY@tc##1{\textcolor[rgb]{0.00,0.00,1.00}{##1}}}
\@namedef{PY@tok@nn}{\let\PY@bf=\textbf\def\PY@tc##1{\textcolor[rgb]{0.00,0.00,1.00}{##1}}}
\@namedef{PY@tok@ne}{\let\PY@bf=\textbf\def\PY@tc##1{\textcolor[rgb]{0.80,0.25,0.22}{##1}}}
\@namedef{PY@tok@nv}{\def\PY@tc##1{\textcolor[rgb]{0.10,0.09,0.49}{##1}}}
\@namedef{PY@tok@no}{\def\PY@tc##1{\textcolor[rgb]{0.53,0.00,0.00}{##1}}}
\@namedef{PY@tok@nl}{\def\PY@tc##1{\textcolor[rgb]{0.46,0.46,0.00}{##1}}}
\@namedef{PY@tok@ni}{\let\PY@bf=\textbf\def\PY@tc##1{\textcolor[rgb]{0.44,0.44,0.44}{##1}}}
\@namedef{PY@tok@na}{\def\PY@tc##1{\textcolor[rgb]{0.41,0.47,0.13}{##1}}}
\@namedef{PY@tok@nt}{\let\PY@bf=\textbf\def\PY@tc##1{\textcolor[rgb]{0.00,0.50,0.00}{##1}}}
\@namedef{PY@tok@nd}{\def\PY@tc##1{\textcolor[rgb]{0.67,0.13,1.00}{##1}}}
\@namedef{PY@tok@s}{\def\PY@tc##1{\textcolor[rgb]{0.73,0.13,0.13}{##1}}}
\@namedef{PY@tok@sd}{\let\PY@it=\textit\def\PY@tc##1{\textcolor[rgb]{0.73,0.13,0.13}{##1}}}
\@namedef{PY@tok@si}{\let\PY@bf=\textbf\def\PY@tc##1{\textcolor[rgb]{0.64,0.35,0.47}{##1}}}
\@namedef{PY@tok@se}{\let\PY@bf=\textbf\def\PY@tc##1{\textcolor[rgb]{0.67,0.36,0.12}{##1}}}
\@namedef{PY@tok@sr}{\def\PY@tc##1{\textcolor[rgb]{0.64,0.35,0.47}{##1}}}
\@namedef{PY@tok@ss}{\def\PY@tc##1{\textcolor[rgb]{0.10,0.09,0.49}{##1}}}
\@namedef{PY@tok@sx}{\def\PY@tc##1{\textcolor[rgb]{0.00,0.50,0.00}{##1}}}
\@namedef{PY@tok@m}{\def\PY@tc##1{\textcolor[rgb]{0.40,0.40,0.40}{##1}}}
\@namedef{PY@tok@gh}{\let\PY@bf=\textbf\def\PY@tc##1{\textcolor[rgb]{0.00,0.00,0.50}{##1}}}
\@namedef{PY@tok@gu}{\let\PY@bf=\textbf\def\PY@tc##1{\textcolor[rgb]{0.50,0.00,0.50}{##1}}}
\@namedef{PY@tok@gd}{\def\PY@tc##1{\textcolor[rgb]{0.63,0.00,0.00}{##1}}}
\@namedef{PY@tok@gi}{\def\PY@tc##1{\textcolor[rgb]{0.00,0.52,0.00}{##1}}}
\@namedef{PY@tok@gr}{\def\PY@tc##1{\textcolor[rgb]{0.89,0.00,0.00}{##1}}}
\@namedef{PY@tok@ge}{\let\PY@it=\textit}
\@namedef{PY@tok@gs}{\let\PY@bf=\textbf}
\@namedef{PY@tok@gp}{\let\PY@bf=\textbf\def\PY@tc##1{\textcolor[rgb]{0.00,0.00,0.50}{##1}}}
\@namedef{PY@tok@go}{\def\PY@tc##1{\textcolor[rgb]{0.44,0.44,0.44}{##1}}}
\@namedef{PY@tok@gt}{\def\PY@tc##1{\textcolor[rgb]{0.00,0.27,0.87}{##1}}}
\@namedef{PY@tok@err}{\def\PY@bc##1{{\setlength{\fboxsep}{\string -\fboxrule}\fcolorbox[rgb]{1.00,0.00,0.00}{1,1,1}{\strut ##1}}}}
\@namedef{PY@tok@kc}{\let\PY@bf=\textbf\def\PY@tc##1{\textcolor[rgb]{0.00,0.50,0.00}{##1}}}
\@namedef{PY@tok@kd}{\let\PY@bf=\textbf\def\PY@tc##1{\textcolor[rgb]{0.00,0.50,0.00}{##1}}}
\@namedef{PY@tok@kn}{\let\PY@bf=\textbf\def\PY@tc##1{\textcolor[rgb]{0.00,0.50,0.00}{##1}}}
\@namedef{PY@tok@kr}{\let\PY@bf=\textbf\def\PY@tc##1{\textcolor[rgb]{0.00,0.50,0.00}{##1}}}
\@namedef{PY@tok@bp}{\def\PY@tc##1{\textcolor[rgb]{0.00,0.50,0.00}{##1}}}
\@namedef{PY@tok@fm}{\def\PY@tc##1{\textcolor[rgb]{0.00,0.00,1.00}{##1}}}
\@namedef{PY@tok@vc}{\def\PY@tc##1{\textcolor[rgb]{0.10,0.09,0.49}{##1}}}
\@namedef{PY@tok@vg}{\def\PY@tc##1{\textcolor[rgb]{0.10,0.09,0.49}{##1}}}
\@namedef{PY@tok@vi}{\def\PY@tc##1{\textcolor[rgb]{0.10,0.09,0.49}{##1}}}
\@namedef{PY@tok@vm}{\def\PY@tc##1{\textcolor[rgb]{0.10,0.09,0.49}{##1}}}
\@namedef{PY@tok@sa}{\def\PY@tc##1{\textcolor[rgb]{0.73,0.13,0.13}{##1}}}
\@namedef{PY@tok@sb}{\def\PY@tc##1{\textcolor[rgb]{0.73,0.13,0.13}{##1}}}
\@namedef{PY@tok@sc}{\def\PY@tc##1{\textcolor[rgb]{0.73,0.13,0.13}{##1}}}
\@namedef{PY@tok@dl}{\def\PY@tc##1{\textcolor[rgb]{0.73,0.13,0.13}{##1}}}
\@namedef{PY@tok@s2}{\def\PY@tc##1{\textcolor[rgb]{0.73,0.13,0.13}{##1}}}
\@namedef{PY@tok@sh}{\def\PY@tc##1{\textcolor[rgb]{0.73,0.13,0.13}{##1}}}
\@namedef{PY@tok@s1}{\def\PY@tc##1{\textcolor[rgb]{0.73,0.13,0.13}{##1}}}
\@namedef{PY@tok@mb}{\def\PY@tc##1{\textcolor[rgb]{0.40,0.40,0.40}{##1}}}
\@namedef{PY@tok@mf}{\def\PY@tc##1{\textcolor[rgb]{0.40,0.40,0.40}{##1}}}
\@namedef{PY@tok@mh}{\def\PY@tc##1{\textcolor[rgb]{0.40,0.40,0.40}{##1}}}
\@namedef{PY@tok@mi}{\def\PY@tc##1{\textcolor[rgb]{0.40,0.40,0.40}{##1}}}
\@namedef{PY@tok@il}{\def\PY@tc##1{\textcolor[rgb]{0.40,0.40,0.40}{##1}}}
\@namedef{PY@tok@mo}{\def\PY@tc##1{\textcolor[rgb]{0.40,0.40,0.40}{##1}}}
\@namedef{PY@tok@ch}{\let\PY@it=\textit\def\PY@tc##1{\textcolor[rgb]{0.24,0.48,0.48}{##1}}}
\@namedef{PY@tok@cm}{\let\PY@it=\textit\def\PY@tc##1{\textcolor[rgb]{0.24,0.48,0.48}{##1}}}
\@namedef{PY@tok@cpf}{\let\PY@it=\textit\def\PY@tc##1{\textcolor[rgb]{0.24,0.48,0.48}{##1}}}
\@namedef{PY@tok@c1}{\let\PY@it=\textit\def\PY@tc##1{\textcolor[rgb]{0.24,0.48,0.48}{##1}}}
\@namedef{PY@tok@cs}{\let\PY@it=\textit\def\PY@tc##1{\textcolor[rgb]{0.24,0.48,0.48}{##1}}}

\def\PYZbs{\char`\\}
\def\PYZus{\char`\_}
\def\PYZob{\char`\{}
\def\PYZcb{\char`\}}
\def\PYZca{\char`\^}
\def\PYZam{\char`\&}
\def\PYZlt{\char`\<}
\def\PYZgt{\char`\>}
\def\PYZsh{\char`\#}
\def\PYZpc{\char`\%}
\def\PYZdl{\char`\$}
\def\PYZhy{\char`\-}
\def\PYZsq{\char`\'}
\def\PYZdq{\char`\"}
\def\PYZti{\char`\~}
% for compatibility with earlier versions
\def\PYZat{@}
\def\PYZlb{[}
\def\PYZrb{]}
\makeatother


    % For linebreaks inside Verbatim environment from package fancyvrb.
    \makeatletter
        \newbox\Wrappedcontinuationbox
        \newbox\Wrappedvisiblespacebox
        \newcommand*\Wrappedvisiblespace {\textcolor{red}{\textvisiblespace}}
        \newcommand*\Wrappedcontinuationsymbol {\textcolor{red}{\llap{\tiny$\m@th\hookrightarrow$}}}
        \newcommand*\Wrappedcontinuationindent {3ex }
        \newcommand*\Wrappedafterbreak {\kern\Wrappedcontinuationindent\copy\Wrappedcontinuationbox}
        % Take advantage of the already applied Pygments mark-up to insert
        % potential linebreaks for TeX processing.
        %        {, <, #, %, $, ' and ": go to next line.
        %        _, }, ^, &, >, - and ~: stay at end of broken line.
        % Use of \textquotesingle for straight quote.
        \newcommand*\Wrappedbreaksatspecials {%
            \def\PYGZus{\discretionary{\char`\_}{\Wrappedafterbreak}{\char`\_}}%
            \def\PYGZob{\discretionary{}{\Wrappedafterbreak\char`\{}{\char`\{}}%
            \def\PYGZcb{\discretionary{\char`\}}{\Wrappedafterbreak}{\char`\}}}%
            \def\PYGZca{\discretionary{\char`\^}{\Wrappedafterbreak}{\char`\^}}%
            \def\PYGZam{\discretionary{\char`\&}{\Wrappedafterbreak}{\char`\&}}%
            \def\PYGZlt{\discretionary{}{\Wrappedafterbreak\char`\<}{\char`\<}}%
            \def\PYGZgt{\discretionary{\char`\>}{\Wrappedafterbreak}{\char`\>}}%
            \def\PYGZsh{\discretionary{}{\Wrappedafterbreak\char`\#}{\char`\#}}%
            \def\PYGZpc{\discretionary{}{\Wrappedafterbreak\char`\%}{\char`\%}}%
            \def\PYGZdl{\discretionary{}{\Wrappedafterbreak\char`\$}{\char`\$}}%
            \def\PYGZhy{\discretionary{\char`\-}{\Wrappedafterbreak}{\char`\-}}%
            \def\PYGZsq{\discretionary{}{\Wrappedafterbreak\textquotesingle}{\textquotesingle}}%
            \def\PYGZdq{\discretionary{}{\Wrappedafterbreak\char`\"}{\char`\"}}%
            \def\PYGZti{\discretionary{\char`\~}{\Wrappedafterbreak}{\char`\~}}%
        }
        % Some characters . , ; ? ! / are not pygmentized.
        % This macro makes them "active" and they will insert potential linebreaks
        \newcommand*\Wrappedbreaksatpunct {%
            \lccode`\~`\.\lowercase{\def~}{\discretionary{\hbox{\char`\.}}{\Wrappedafterbreak}{\hbox{\char`\.}}}%
            \lccode`\~`\,\lowercase{\def~}{\discretionary{\hbox{\char`\,}}{\Wrappedafterbreak}{\hbox{\char`\,}}}%
            \lccode`\~`\;\lowercase{\def~}{\discretionary{\hbox{\char`\;}}{\Wrappedafterbreak}{\hbox{\char`\;}}}%
            \lccode`\~`\:\lowercase{\def~}{\discretionary{\hbox{\char`\:}}{\Wrappedafterbreak}{\hbox{\char`\:}}}%
            \lccode`\~`\?\lowercase{\def~}{\discretionary{\hbox{\char`\?}}{\Wrappedafterbreak}{\hbox{\char`\?}}}%
            \lccode`\~`\!\lowercase{\def~}{\discretionary{\hbox{\char`\!}}{\Wrappedafterbreak}{\hbox{\char`\!}}}%
            \lccode`\~`\/\lowercase{\def~}{\discretionary{\hbox{\char`\/}}{\Wrappedafterbreak}{\hbox{\char`\/}}}%
            \catcode`\.\active
            \catcode`\,\active
            \catcode`\;\active
            \catcode`\:\active
            \catcode`\?\active
            \catcode`\!\active
            \catcode`\/\active
            \lccode`\~`\~
        }
    \makeatother

    \let\OriginalVerbatim=\Verbatim
    \makeatletter
    \renewcommand{\Verbatim}[1][1]{%
        %\parskip\z@skip
        \sbox\Wrappedcontinuationbox {\Wrappedcontinuationsymbol}%
        \sbox\Wrappedvisiblespacebox {\FV@SetupFont\Wrappedvisiblespace}%
        \def\FancyVerbFormatLine ##1{\hsize\linewidth
            \vtop{\raggedright\hyphenpenalty\z@\exhyphenpenalty\z@
                \doublehyphendemerits\z@\finalhyphendemerits\z@
                \strut ##1\strut}%
        }%
        % If the linebreak is at a space, the latter will be displayed as visible
        % space at end of first line, and a continuation symbol starts next line.
        % Stretch/shrink are however usually zero for typewriter font.
        \def\FV@Space {%
            \nobreak\hskip\z@ plus\fontdimen3\font minus\fontdimen4\font
            \discretionary{\copy\Wrappedvisiblespacebox}{\Wrappedafterbreak}
            {\kern\fontdimen2\font}%
        }%

        % Allow breaks at special characters using \PYG... macros.
        \Wrappedbreaksatspecials
        % Breaks at punctuation characters . , ; ? ! and / need catcode=\active
        \OriginalVerbatim[#1,codes*=\Wrappedbreaksatpunct]%
    }
    \makeatother

    % Exact colors from NB
    \definecolor{incolor}{HTML}{303F9F}
    \definecolor{outcolor}{HTML}{D84315}
    \definecolor{cellborder}{HTML}{CFCFCF}
    \definecolor{cellbackground}{HTML}{F7F7F7}

    % prompt
    \makeatletter
    \newcommand{\boxspacing}{\kern\kvtcb@left@rule\kern\kvtcb@boxsep}
    \makeatother
    \newcommand{\prompt}[4]{
        {\ttfamily\llap{{\color{#2}[#3]:\hspace{3pt}#4}}\vspace{-\baselineskip}}
    }
    

    
    % Prevent overflowing lines due to hard-to-break entities
    \sloppy
    % Setup hyperref package
    \hypersetup{
      breaklinks=true,  % so long urls are correctly broken across lines
      colorlinks=true,
      urlcolor=urlcolor,
      linkcolor=linkcolor,
      citecolor=citecolor,
      }
    % Slightly bigger margins than the latex defaults
    
    \geometry{verbose,tmargin=1in,bmargin=1in,lmargin=1in,rmargin=1in}
    
    

\begin{document}
    
    \maketitle
    
    

    
    \hypertarget{introduction}{%
\section{Introduction}\label{introduction}}

    \textbf{File:} Notebook\_BVP\_to\_end.ipynb\\
\textbf{Author:} Elizabeth Gould\\
\textbf{Date:} 11.03.2026\\
\textbf{Problem:} Find the dependence of the current of a loop on the
parameters \(B\), \(k\), \(R\), and \(\mu\).

    This Jupyter notebook is used for modeling the relative persistent
current \(I\) in micrometer scale gold loops which solve the BVP of the
nonlinear Schrödinger equation.

For our observed approximate solution to the nonlinear Schrödinger
equation \begin{equation}
\psi(x) = e^{i M_{j} x} \left( a e^{i k_{j} x} + b e^{- i k_{j} x} \right),
\end{equation} we know that \(\psi(0) = \psi(2\pi R)\) when
\begin{equation}
\psi^\prime(0) = i M_{j} + k_j \left( \frac{e^{- 2 i M_{j}\pi R} - \cos\left( 2\pi k_{j} R \right)}{\sin\left(2\pi k_{j} R\right)}\right),
\end{equation}

where \(k_{j}\left(\mu, R, B, k_0, \psi^\prime(0) \right)\) and
\(M_{j}\left(\mu, R, B, k_0, \psi^\prime(0) \right)\) are given by our
previously solved model. Since the equations for \(k_j\) and \(M_j\) of
our model are dependent on \(\psi^\prime(0)\), we have obtained the
values of \(\psi^\prime(0)\) which solve this BVP equation with a
numerical root finding algorithm. We now wish to take the task to its
conclusion, and determine how the altered solutions change our original
solution to the relative current between two rings in the system.

    As will be seen, and as should be obvious from our BVP matching
condition, the solution to this problem is the same as for the linear
case,

\begin{equation}
I_j = - \left| \psi(0) \right|^2 \frac{\sin\left(2 \pi M_j R_j\right)}{\sin\left(2 \pi k_j R_j\right)},
\end{equation}

where \(\psi(0)\) is our value of \(\psi(x)\) at our point of
intersection. This means that our final result for the ratio of current
between loops is given by \begin{equation}
\frac{I_j}{I_n} = \frac{\sin\left(2 \pi M_j R_j\right)\sin\left(2 \pi k_n R_n\right)}{\sin\left(2 \pi k_j R_j\right)\sin\left(2 \pi M_n R_n\right)},
\end{equation} as was the case for the linear case. Note that our values
of \(\mu\) in our expressions for
\(k_j\left(\mu, R_j, B, k_0, \psi^\prime_j(0) \right)\) and
\(M_j\left(\mu, R_j, B, k_0, \psi^\prime_j(0) \right)\) assume
\(\psi(0) = 1\), and for the case when \(\psi(0) \neq 1\), the \(\mu\)
in our results is given by
\(\mu_{eff} = \mu_{A_0} \left| \psi(0) \right|^2\), where \(\mu_{A_0}\)
is the value of \(\mu\) which was believed have been provided to the
Schrödinger equation.

However, there are a few complications. First, \(k_j\) and \(M_j\) are
now given by a function which depends on \(\psi^\prime(0)\). This means
we need a new model for \(k_j\) and \(M_j\) which solve the BVP rather
than which solve the IVP. The second complication is that we don't know
\(\mu\). More specifically, we don't know how to match \(\mu\) between
rings -- is \(\mu\) effectively equal for equal initial wave function,
or for equal maximum wave function.

I will proceed by assuming equal \(\mu\) for equal initial wave
function. This is not because this is the result I believe to be the
case, but because it is easier, and therefore useful for an initial
analysis. As can be seen from our equation for \(\mu\), \begin{equation}
\mu_{eff} = \mu_{A_0} \left| \psi(0) \right|^2,
\end{equation} if \(\psi(0)\) and \(\mu_{A_0}\) are both equal for all
rings in the system, then \(\mu_{eff}\) is also equal for all of the
rings, and the solution for relative current can easily be analyzed by
just choosing a value for \(\mu_{A_0}\) or \(\mu_{eff}\).

If we are to believe that instead \(\mu\) needs to be matched for equal
\(\left|\psi\right|_{max}\), our value of \(\mu_{eff}\) is now loop
dependent, and given by \begin{equation}
\mu_{eff} = \mu_{max} \left| \psi \right|_{max}^2.
\end{equation} For a fixed \(\mu_{max}\), we need to calculate
\(\mu_{eff, j}\) for every ring by numerically finding the following
roots, \begin{equation}
0 = \mu_{eff} - \mu_{max} \left| \psi \right| _{max}^2 = \mu_{eff, j} - \mu_{max} \left(\left| a_j \right|^2 + \left| b_j \right|^2 + 2 \left| a_j \right|\left| b_j \right|\right)
\end{equation} where \begin{equation}
a_j = \psi(0) \frac{1 - e^{i \left(-k_j + M_j \right) 2 \pi R_j}}{e^{i 2 M_j \pi R_j} 2 i \sin \left( 2 \pi k_j R_j \right)},
\end{equation}

\begin{equation}
b_j = \psi(0) \frac{-1 + e^{i \left(k_j + M_j \right) 2 \pi R_j}}{e^{i 2 M_j \pi R_j} 2 i \sin \left( 2 \pi k_j R_j \right)},
\end{equation} and \(k_j \left(\mu_{eff,j}, R_j, B, k_0\right)\) and
\(M_j \left(\mu_{eff,j}, R_j, B, k_0 \right)\) are the values of
\(k_j\left(\mu_{eff,j}, R_j, B, k_0, \psi^\prime_j(0) \right)\) and
\(M_j\left(\mu_{eff,j}, R_j, B, k_0, \psi^\prime_j(0) \right)\) which
solve the BVP.

    After demonstrating that our solution is indeed the case, I will, in my
next and final notebook, go back to modelling and analyze the
dependences of \(k_j\left(\mu, R, B, k_0 \right)\) and
\(M_j\left(\mu, R, B, k_0 \right)\) on the parameters \(R\), \(B\),
\(dk\), and \(\mu\), with \(k_0\) given by \(k_0 = k_{Fermi, Au} + dk\).

    For completeness, the final boundary matching term is \begin{equation}
i k_0 \frac{r - 1}{r+1} = i \sum^{n}_{j = 1}k_j\left( a_j \left( 1 - e^{i \left(k_j + M_j\right) 2 \pi R_j} \right) - b_j \left( 1 - e^{i \left(-k_j + M_j\right) 2 \pi R_j}\right) \right).
\end{equation}

Since \(\psi(0) = 1 + r\) in the linear analysis, but \(\psi(0) = 1\)
for our analysis, we need to divide by \(1+r\). Otherwise the equation
is the same, just for new values for \(k_j\) and \(M_j\). This equation
is solved for \(r\), which is our reflection coefficient. The analysis
is in theory the same as for the linear case, but our right hand side,
which is independent of \(r\), has a new value. We effectively need to
solve the equation \begin{equation}
\frac{r - 1}{r+1} = W,
\end{equation} where \(W\) is a function of \(n\), \(k_0\), \(B\),
\(\mu\) and all of our \(R_j\) values. This is the same as for the
linear case, we have just changed our value \(W\). Note, that like the
linear case, the complete system, and thus \(n\) as well, are
unimportant for the analysis of our relative currents.

    \hypertarget{load-grids}{%
\section{Load Grids}\label{load-grids}}

    We will start, as typical, with our libraries.

    \begin{tcolorbox}[breakable, size=fbox, boxrule=1pt, pad at break*=1mm,colback=cellbackground, colframe=cellborder]
\prompt{In}{incolor}{ }{\boxspacing}
\begin{Verbatim}[commandchars=\\\{\}]
\PY{c+c1}{\PYZsh{} my library}
\PY{k+kn}{import} \PY{n+nn}{eelib}
\PY{k+kn}{from} \PY{n+nn}{eelib} \PY{k+kn}{import} \PY{n}{pi}\PY{p}{,} \PY{n}{B\PYZus{}max}\PY{p}{,} \PY{n}{phi0inv}\PY{p}{,} \PY{n}{R\PYZus{}max}
\PY{k+kn}{from} \PY{n+nn}{eelib} \PY{k+kn}{import} \PY{n}{clean\PYZus{}table}\PY{p}{,} \PY{n}{clean\PYZus{}table\PYZus{}MC}

\PY{k+kn}{import} \PY{n+nn}{numpy} \PY{k}{as} \PY{n+nn}{np}
\PY{k+kn}{import} \PY{n+nn}{pickle}
\PY{k+kn}{import} \PY{n+nn}{pandas} \PY{k}{as} \PY{n+nn}{pd}
\PY{k+kn}{import} \PY{n+nn}{warnings}

\PY{c+c1}{\PYZsh{} machine learning}
\PY{k+kn}{from} \PY{n+nn}{sklearn}\PY{n+nn}{.}\PY{n+nn}{linear\PYZus{}model} \PY{k+kn}{import} \PY{n}{LinearRegression}
\PY{k+kn}{from} \PY{n+nn}{sklearn}\PY{n+nn}{.}\PY{n+nn}{metrics} \PY{k+kn}{import} \PY{n}{r2\PYZus{}score}\PY{p}{,} \PY{n}{mean\PYZus{}squared\PYZus{}error}\PY{p}{,} \PY{n}{mean\PYZus{}absolute\PYZus{}error}

\PY{c+c1}{\PYZsh{} plotting}
\PY{k+kn}{import} \PY{n+nn}{matplotlib}\PY{n+nn}{.}\PY{n+nn}{pyplot} \PY{k}{as} \PY{n+nn}{plt}
\PY{k+kn}{import} \PY{n+nn}{seaborn} \PY{k}{as} \PY{n+nn}{sns} \PY{c+c1}{\PYZsh{}statistical data visualization}
\end{Verbatim}
\end{tcolorbox}

    And loading our data from previous runs of our bvp\_grid code.

    \begin{tcolorbox}[breakable, size=fbox, boxrule=1pt, pad at break*=1mm,colback=cellbackground, colframe=cellborder]
\prompt{In}{incolor}{2}{\boxspacing}
\begin{Verbatim}[commandchars=\\\{\}]
\PY{n}{filename} \PY{o}{=} \PY{l+s+s1}{\PYZsq{}}\PY{l+s+s1}{grid020}\PY{l+s+s1}{\PYZsq{}}

\PY{c+c1}{\PYZsh{} loading script; be careful of python version}
\PY{n}{file} \PY{o}{=} \PY{n+nb}{open}\PY{p}{(}\PY{n}{filename}\PY{p}{,} \PY{l+s+s1}{\PYZsq{}}\PY{l+s+s1}{rb}\PY{l+s+s1}{\PYZsq{}}\PY{p}{)}    
\PY{n}{gridl\PYZus{}1} \PY{o}{=} \PY{n}{pickle}\PY{o}{.}\PY{n}{load}\PY{p}{(}\PY{n}{file}\PY{p}{)}
\PY{n}{file}\PY{o}{.}\PY{n}{close}\PY{p}{(}\PY{p}{)}
\end{Verbatim}
\end{tcolorbox}
\begin{tcolorbox}[breakable, size=fbox, boxrule=1pt, pad at break*=1mm,colback=cellbackground, colframe=cellborder]
\prompt{In}{incolor}{41}{\boxspacing}
\begin{Verbatim}[commandchars=\\\{\}]
filename = 'grid021'

# loading script; be careful of python version
file = open(filename, 'rb')    
gridl_2 = pickle.load(file)
file.close()
\end{Verbatim}
\end{tcolorbox}
    \begin{tcolorbox}[breakable, size=fbox, boxrule=1pt, pad at break*=1mm,colback=cellbackground, colframe=cellborder]
\prompt{In}{incolor}{3}{\boxspacing}
\begin{Verbatim}[commandchars=\\\{\}]
\PY{n}{filename} \PY{o}{=} \PY{l+s+s1}{\PYZsq{}}\PY{l+s+s1}{grid022}\PY{l+s+s1}{\PYZsq{}}

\PY{c+c1}{\PYZsh{} loading script; be careful of python version}
\PY{n}{file} \PY{o}{=} \PY{n+nb}{open}\PY{p}{(}\PY{n}{filename}\PY{p}{,} \PY{l+s+s1}{\PYZsq{}}\PY{l+s+s1}{rb}\PY{l+s+s1}{\PYZsq{}}\PY{p}{)}    
\PY{n}{gridl\PYZus{}3} \PY{o}{=} \PY{n}{pickle}\PY{o}{.}\PY{n}{load}\PY{p}{(}\PY{n}{file}\PY{p}{)}
\PY{n}{file}\PY{o}{.}\PY{n}{close}\PY{p}{(}\PY{p}{)}
\end{Verbatim}
\end{tcolorbox}
\begin{tcolorbox}[breakable, size=fbox, boxrule=1pt, pad at break*=1mm,colback=cellbackground, colframe=cellborder]
\prompt{In}{incolor}{41}{\boxspacing}
\begin{Verbatim}[commandchars=\\\{\}]
filename = 'grid023'

# loading script; be careful of python version
file = open(filename, 'rb')    
gridl_4 = pickle.load(file)
file.close()
\end{Verbatim}
\end{tcolorbox}
    \begin{tcolorbox}[breakable, size=fbox, boxrule=1pt, pad at break*=1mm,colback=cellbackground, colframe=cellborder]
\prompt{In}{incolor}{4}{\boxspacing}
\begin{Verbatim}[commandchars=\\\{\}]
\PY{n}{filename} \PY{o}{=} \PY{l+s+s1}{\PYZsq{}}\PY{l+s+s1}{grid025}\PY{l+s+s1}{\PYZsq{}}

\PY{c+c1}{\PYZsh{} loading script; be careful of python version}
\PY{n}{file} \PY{o}{=} \PY{n+nb}{open}\PY{p}{(}\PY{n}{filename}\PY{p}{,} \PY{l+s+s1}{\PYZsq{}}\PY{l+s+s1}{rb}\PY{l+s+s1}{\PYZsq{}}\PY{p}{)}    
\PY{n}{gridl\PYZus{}5} \PY{o}{=} \PY{n}{pickle}\PY{o}{.}\PY{n}{load}\PY{p}{(}\PY{n}{file}\PY{p}{)}
\PY{n}{file}\PY{o}{.}\PY{n}{close}\PY{p}{(}\PY{p}{)}
\end{Verbatim}
\end{tcolorbox}

    \begin{tcolorbox}[breakable, size=fbox, boxrule=1pt, pad at break*=1mm,colback=cellbackground, colframe=cellborder]
\prompt{In}{incolor}{5}{\boxspacing}
\begin{Verbatim}[commandchars=\\\{\}]
\PY{n}{filename} \PY{o}{=} \PY{l+s+s1}{\PYZsq{}}\PY{l+s+s1}{grid024}\PY{l+s+s1}{\PYZsq{}}

\PY{c+c1}{\PYZsh{} loading script; be careful of python version}
\PY{n}{file} \PY{o}{=} \PY{n+nb}{open}\PY{p}{(}\PY{n}{filename}\PY{p}{,} \PY{l+s+s1}{\PYZsq{}}\PY{l+s+s1}{rb}\PY{l+s+s1}{\PYZsq{}}\PY{p}{)}    
\PY{n}{gridl} \PY{o}{=} \PY{n}{pickle}\PY{o}{.}\PY{n}{load}\PY{p}{(}\PY{n}{file}\PY{p}{)}
\PY{n}{file}\PY{o}{.}\PY{n}{close}\PY{p}{(}\PY{p}{)}
\end{Verbatim}
\end{tcolorbox}

    Now to check our objects.

    \begin{tcolorbox}[breakable, size=fbox, boxrule=1pt, pad at break*=1mm,colback=cellbackground, colframe=cellborder]
\prompt{In}{incolor}{6}{\boxspacing}
\begin{Verbatim}[commandchars=\\\{\}]
\PY{n+nb}{print}\PY{p}{(}\PY{n}{gridl\PYZus{}1}\PY{p}{)}
\PY{n+nb}{print}\PY{p}{(}\PY{l+s+s2}{\PYZdq{}}\PY{l+s+s2}{\PYZdq{}}\PY{p}{)}
\PY{n+nb}{print}\PY{p}{(}\PY{n}{gridl\PYZus{}5}\PY{p}{)}
\end{Verbatim}
\end{tcolorbox}

    \begin{Verbatim}[commandchars=\\\{\}]
Grid object to measure current:
mu has 11 points from 1e-10 to 1e-08.
dk has 11 points from 0.51 to 0.515.
B has 11 points from 0.5 to 0.50055.
R has 11 points from 0.9 to 0.90003.
A is: [1.]
k0 is: 12000000000.0
Total number of parameters: 14641

Monte Carlo object to measure current:
mu has points from 1e-07 to 1e-07.
dk has points from 0.51000044582215 to 0.514996169805716.
B has points from 0.5000003637516297 to 0.500549933737302.
R has points from 0.9000000031221171 to 0.9000299973995092.
A has points from 1.0 to 1.0.
k0 has points from 12000000000.0 to 12000000000.0.
Number of points: 1000
    \end{Verbatim}

    \begin{tcolorbox}[breakable, size=fbox, boxrule=1pt, pad at break*=1mm,colback=cellbackground, colframe=cellborder]
\prompt{In}{incolor}{7}{\boxspacing}
\begin{Verbatim}[commandchars=\\\{\}]
\PY{n+nb}{print}\PY{p}{(}\PY{n}{gridl\PYZus{}3}\PY{p}{)}
\PY{n+nb}{print}\PY{p}{(}\PY{l+s+s2}{\PYZdq{}}\PY{l+s+s2}{\PYZdq{}}\PY{p}{)}
\PY{n+nb}{print}\PY{p}{(}\PY{n}{gridl}\PY{p}{)}
\end{Verbatim}
\end{tcolorbox}

    \begin{Verbatim}[commandchars=\\\{\}]
Grid object to measure current:
mu has 11 points from 1e-10 to 1e-08.
dk has 11 points from 0.1 to 0.9.
B has 11 points from 0.11 to 1.0.
R has 11 points from 0.57 to 1.0.
A is: [1.]
k0 is: 12000000000.0
Total number of parameters: 14641

Grid object to measure current:
mu has 6 points from 1e-10 to 1e-08.
dk has 6 points from 0.51 to 0.515.
B has 6 points from 0.5 to 0.50055.
R has 6 points from 0.9 to 0.90003.
A is: [1.]
k0 is: 12000000000.0
Total number of parameters: 1296
    \end{Verbatim}

    My grids are of the following forms: 1) gridl\_1 is the grid for small
parameter changes, as is gridl. 2) gridl\_2 is the Monte Carlo run for
small parameter changes, as is grid\_5. 3) gridl\_3 is the grid for
large parameter changes. 4) gridl\_4 is the Monte Carlo run for large
parameter changes.

gridl and gridl\_5 are the older, smaller datasets, which is the reason
why there are two runs with similar parameters. However, I had
accidentally run gridl\_2 and gridl\_4 with a very small number of
parameters, so I have removed them from the code. I didn't redo these
runs, as this analysis had since been proven unnecessary.

Small parameter changes have been chosen to analyze the rapid
oscillations, which I know exist in my model. Larger parameter changes
attempt to find the large-scale behavior, like slow oscillations or
linear depencies. I know they will be ineffective at this task without
fully removing the fast oscillations, but I have the data here for
trying to determine a starting point of analysis.

    \hypertarget{make-tables}{%
\section{Make Tables}\label{make-tables}}

    First, load the tables from my objects into pandas dataframes. The data
has been saved in my objects in a form of a list of dictionaries,
designed to easily form a pandas dataframe with a simple function call.

Tables below can be commented or uncommented as needed.

    \begin{tcolorbox}[breakable, size=fbox, boxrule=1pt, pad at break*=1mm,colback=cellbackground, colframe=cellborder]
\prompt{In}{incolor}{9}{\boxspacing}
\begin{Verbatim}[commandchars=\\\{\}]
\PY{n}{tbl}    \PY{o}{=} \PY{n}{pd}\PY{o}{.}\PY{n}{DataFrame}\PY{p}{(}\PY{n}{gridl}\PY{o}{.}\PY{n}{derivs}\PY{p}{)}
\PY{n}{tbl001} \PY{o}{=} \PY{n}{pd}\PY{o}{.}\PY{n}{DataFrame}\PY{p}{(}\PY{n}{gridl\PYZus{}1}\PY{o}{.}\PY{n}{derivs}\PY{p}{)}
\PY{c+c1}{\PYZsh{}tbl002 = pd.DataFrame(gridl\PYZus{}2.derivs)}
\PY{n}{tbl003} \PY{o}{=} \PY{n}{pd}\PY{o}{.}\PY{n}{DataFrame}\PY{p}{(}\PY{n}{gridl\PYZus{}3}\PY{o}{.}\PY{n}{derivs}\PY{p}{)}
\PY{c+c1}{\PYZsh{}tbl004 = pd.DataFrame(gridl\PYZus{}4.derivs)}
\PY{n}{tbl005} \PY{o}{=} \PY{n}{pd}\PY{o}{.}\PY{n}{DataFrame}\PY{p}{(}\PY{n}{gridl\PYZus{}5}\PY{o}{.}\PY{n}{derivs}\PY{p}{)}
\end{Verbatim}
\end{tcolorbox}

    I will also need to get the values of my parameters for grid analyses,
so that I can easily select all points with the same value of the given
parameter. We also will pull \(k = k_{Fermi, Au}\) and
\(A = \psi(0) = 1\).

    \begin{tcolorbox}[breakable, size=fbox, boxrule=1pt, pad at break*=1mm,colback=cellbackground, colframe=cellborder]
\prompt{In}{incolor}{10}{\boxspacing}
\begin{Verbatim}[commandchars=\\\{\}]
\PY{n}{R\PYZus{}arr} \PY{o}{=} \PY{n}{tbl}\PY{p}{[}\PY{l+s+s2}{\PYZdq{}}\PY{l+s+s2}{R}\PY{l+s+s2}{\PYZdq{}}\PY{p}{]}\PY{o}{.}\PY{n}{unique}\PY{p}{(}\PY{p}{)}
\PY{n}{B\PYZus{}arr} \PY{o}{=} \PY{n}{tbl}\PY{p}{[}\PY{l+s+s2}{\PYZdq{}}\PY{l+s+s2}{B}\PY{l+s+s2}{\PYZdq{}}\PY{p}{]}\PY{o}{.}\PY{n}{unique}\PY{p}{(}\PY{p}{)}
\PY{n}{dk\PYZus{}arr} \PY{o}{=} \PY{n}{tbl}\PY{p}{[}\PY{l+s+s2}{\PYZdq{}}\PY{l+s+s2}{dk}\PY{l+s+s2}{\PYZdq{}}\PY{p}{]}\PY{o}{.}\PY{n}{unique}\PY{p}{(}\PY{p}{)}
\PY{n}{mu\PYZus{}arr} \PY{o}{=} \PY{n}{tbl}\PY{p}{[}\PY{l+s+s2}{\PYZdq{}}\PY{l+s+s2}{mu}\PY{l+s+s2}{\PYZdq{}}\PY{p}{]}\PY{o}{.}\PY{n}{unique}\PY{p}{(}\PY{p}{)}
\PY{n}{k\PYZus{}val} \PY{o}{=} \PY{n}{tbl}\PY{p}{[}\PY{l+s+s2}{\PYZdq{}}\PY{l+s+s2}{k}\PY{l+s+s2}{\PYZdq{}}\PY{p}{]}\PY{o}{.}\PY{n}{unique}\PY{p}{(}\PY{p}{)}\PY{p}{[}\PY{l+m+mi}{0}\PY{p}{]}
\PY{n}{A\PYZus{}val} \PY{o}{=} \PY{n}{tbl}\PY{p}{[}\PY{l+s+s2}{\PYZdq{}}\PY{l+s+s2}{A}\PY{l+s+s2}{\PYZdq{}}\PY{p}{]}\PY{o}{.}\PY{n}{unique}\PY{p}{(}\PY{p}{)}\PY{p}{[}\PY{l+m+mi}{0}\PY{p}{]}
\end{Verbatim}
\end{tcolorbox}

    \begin{tcolorbox}[breakable, size=fbox, boxrule=1pt, pad at break*=1mm,colback=cellbackground, colframe=cellborder]
\prompt{In}{incolor}{11}{\boxspacing}
\begin{Verbatim}[commandchars=\\\{\}]
\PY{n}{R3\PYZus{}arr} \PY{o}{=} \PY{n}{tbl001}\PY{p}{[}\PY{l+s+s2}{\PYZdq{}}\PY{l+s+s2}{R}\PY{l+s+s2}{\PYZdq{}}\PY{p}{]}\PY{o}{.}\PY{n}{unique}\PY{p}{(}\PY{p}{)}
\PY{n}{B3\PYZus{}arr} \PY{o}{=} \PY{n}{tbl001}\PY{p}{[}\PY{l+s+s2}{\PYZdq{}}\PY{l+s+s2}{B}\PY{l+s+s2}{\PYZdq{}}\PY{p}{]}\PY{o}{.}\PY{n}{unique}\PY{p}{(}\PY{p}{)}
\PY{n}{dk3\PYZus{}arr} \PY{o}{=} \PY{n}{tbl001}\PY{p}{[}\PY{l+s+s2}{\PYZdq{}}\PY{l+s+s2}{dk}\PY{l+s+s2}{\PYZdq{}}\PY{p}{]}\PY{o}{.}\PY{n}{unique}\PY{p}{(}\PY{p}{)}
\PY{n}{mu3\PYZus{}arr} \PY{o}{=} \PY{n}{tbl001}\PY{p}{[}\PY{l+s+s2}{\PYZdq{}}\PY{l+s+s2}{mu}\PY{l+s+s2}{\PYZdq{}}\PY{p}{]}\PY{o}{.}\PY{n}{unique}\PY{p}{(}\PY{p}{)}
\PY{n}{k\PYZus{}val} \PY{o}{=} \PY{n}{tbl001}\PY{p}{[}\PY{l+s+s2}{\PYZdq{}}\PY{l+s+s2}{k}\PY{l+s+s2}{\PYZdq{}}\PY{p}{]}\PY{o}{.}\PY{n}{unique}\PY{p}{(}\PY{p}{)}\PY{p}{[}\PY{l+m+mi}{0}\PY{p}{]}
\PY{n}{A\PYZus{}val} \PY{o}{=} \PY{n}{tbl001}\PY{p}{[}\PY{l+s+s2}{\PYZdq{}}\PY{l+s+s2}{A}\PY{l+s+s2}{\PYZdq{}}\PY{p}{]}\PY{o}{.}\PY{n}{unique}\PY{p}{(}\PY{p}{)}\PY{p}{[}\PY{l+m+mi}{0}\PY{p}{]}
\end{Verbatim}
\end{tcolorbox}

    \begin{tcolorbox}[breakable, size=fbox, boxrule=1pt, pad at break*=1mm,colback=cellbackground, colframe=cellborder]
\prompt{In}{incolor}{12}{\boxspacing}
\begin{Verbatim}[commandchars=\\\{\}]
\PY{n}{R2\PYZus{}arr} \PY{o}{=} \PY{n}{tbl003}\PY{p}{[}\PY{l+s+s2}{\PYZdq{}}\PY{l+s+s2}{R}\PY{l+s+s2}{\PYZdq{}}\PY{p}{]}\PY{o}{.}\PY{n}{unique}\PY{p}{(}\PY{p}{)}
\PY{n}{B2\PYZus{}arr} \PY{o}{=} \PY{n}{tbl003}\PY{p}{[}\PY{l+s+s2}{\PYZdq{}}\PY{l+s+s2}{B}\PY{l+s+s2}{\PYZdq{}}\PY{p}{]}\PY{o}{.}\PY{n}{unique}\PY{p}{(}\PY{p}{)}
\PY{n}{dk2\PYZus{}arr} \PY{o}{=} \PY{n}{tbl003}\PY{p}{[}\PY{l+s+s2}{\PYZdq{}}\PY{l+s+s2}{dk}\PY{l+s+s2}{\PYZdq{}}\PY{p}{]}\PY{o}{.}\PY{n}{unique}\PY{p}{(}\PY{p}{)}
\PY{n}{mu2\PYZus{}arr} \PY{o}{=} \PY{n}{tbl003}\PY{p}{[}\PY{l+s+s2}{\PYZdq{}}\PY{l+s+s2}{mu}\PY{l+s+s2}{\PYZdq{}}\PY{p}{]}\PY{o}{.}\PY{n}{unique}\PY{p}{(}\PY{p}{)}
\PY{n}{k\PYZus{}val} \PY{o}{=} \PY{n}{tbl003}\PY{p}{[}\PY{l+s+s2}{\PYZdq{}}\PY{l+s+s2}{k}\PY{l+s+s2}{\PYZdq{}}\PY{p}{]}\PY{o}{.}\PY{n}{unique}\PY{p}{(}\PY{p}{)}\PY{p}{[}\PY{l+m+mi}{0}\PY{p}{]}
\PY{n}{A\PYZus{}val} \PY{o}{=} \PY{n}{tbl003}\PY{p}{[}\PY{l+s+s2}{\PYZdq{}}\PY{l+s+s2}{A}\PY{l+s+s2}{\PYZdq{}}\PY{p}{]}\PY{o}{.}\PY{n}{unique}\PY{p}{(}\PY{p}{)}\PY{p}{[}\PY{l+m+mi}{0}\PY{p}{]}
\end{Verbatim}
\end{tcolorbox}

    Lets now check to see that the solutions we are given produce the
correct final value of \(\psi(2\pi R)\) according to our models for
\(k_j\) and \(M_j\). The values for ``New End'' and ``Exact End'' on our
table are the values for \(\psi(2\pi R)\) for our nonlinear and linear
solutions, respectively. The have been calculated from our model of
\(\psi(x)\), for said solution. They have not been calculated with the
IVP solver, as the error would make such results meaningless.

    \begin{tcolorbox}[breakable, size=fbox, boxrule=1pt, pad at break*=1mm,colback=cellbackground, colframe=cellborder]
\prompt{In}{incolor}{13}{\boxspacing}
\begin{Verbatim}[commandchars=\\\{\}]
\PY{n+nb}{print}\PY{p}{(}\PY{l+s+s2}{\PYZdq{}}\PY{l+s+s2}{(Table 0) Maximum error with ee interaction:   }\PY{l+s+s2}{\PYZdq{}}\PY{p}{,} \PY{n}{np}\PY{o}{.}\PY{n}{max}\PY{p}{(}\PY{n}{np}\PY{o}{.}\PY{n}{abs}\PY{p}{(}\PY{n}{tbl}\PY{p}{[}\PY{l+s+s2}{\PYZdq{}}\PY{l+s+s2}{New End}\PY{l+s+s2}{\PYZdq{}}\PY{p}{]}\PY{o}{.}\PY{n}{to\PYZus{}numpy}\PY{p}{(}\PY{p}{)}\PY{o}{\PYZhy{}}\PY{p}{(}\PY{l+m+mf}{1.0}\PY{o}{+}\PY{l+m+mf}{0.0}\PY{n}{j}\PY{p}{)}\PY{p}{)}\PY{p}{)}\PY{p}{)}
\PY{n+nb}{print}\PY{p}{(}\PY{l+s+s2}{\PYZdq{}}\PY{l+s+s2}{(Table 0) Maximum error without ee interaction:}\PY{l+s+s2}{\PYZdq{}}\PY{p}{,} \PY{n}{np}\PY{o}{.}\PY{n}{max}\PY{p}{(}\PY{n}{np}\PY{o}{.}\PY{n}{abs}\PY{p}{(}\PY{n}{tbl}\PY{p}{[}\PY{l+s+s2}{\PYZdq{}}\PY{l+s+s2}{Exact End}\PY{l+s+s2}{\PYZdq{}}\PY{p}{]}\PY{o}{.}\PY{n}{to\PYZus{}numpy}\PY{p}{(}\PY{p}{)}\PY{o}{\PYZhy{}}\PY{p}{(}\PY{l+m+mf}{1.0}\PY{o}{+}\PY{l+m+mf}{0.0}\PY{n}{j}\PY{p}{)}\PY{p}{)}\PY{p}{)}\PY{p}{)}
\PY{n+nb}{print}\PY{p}{(}\PY{l+s+s2}{\PYZdq{}}\PY{l+s+s2}{(Table 1) Maximum error with ee interaction:   }\PY{l+s+s2}{\PYZdq{}}\PY{p}{,} \PY{n}{np}\PY{o}{.}\PY{n}{max}\PY{p}{(}\PY{n}{np}\PY{o}{.}\PY{n}{abs}\PY{p}{(}\PY{n}{tbl001}\PY{p}{[}\PY{l+s+s2}{\PYZdq{}}\PY{l+s+s2}{New End}\PY{l+s+s2}{\PYZdq{}}\PY{p}{]}\PY{o}{.}\PY{n}{to\PYZus{}numpy}\PY{p}{(}\PY{p}{)}\PY{o}{\PYZhy{}}\PY{p}{(}\PY{l+m+mf}{1.0}\PY{o}{+}\PY{l+m+mf}{0.0}\PY{n}{j}\PY{p}{)}\PY{p}{)}\PY{p}{)}\PY{p}{)}
\PY{n+nb}{print}\PY{p}{(}\PY{l+s+s2}{\PYZdq{}}\PY{l+s+s2}{(Table 1) Maximum error without ee interaction:}\PY{l+s+s2}{\PYZdq{}}\PY{p}{,} \PY{n}{np}\PY{o}{.}\PY{n}{max}\PY{p}{(}\PY{n}{np}\PY{o}{.}\PY{n}{abs}\PY{p}{(}\PY{n}{tbl001}\PY{p}{[}\PY{l+s+s2}{\PYZdq{}}\PY{l+s+s2}{Exact End}\PY{l+s+s2}{\PYZdq{}}\PY{p}{]}\PY{o}{.}\PY{n}{to\PYZus{}numpy}\PY{p}{(}\PY{p}{)}\PY{o}{\PYZhy{}}\PY{p}{(}\PY{l+m+mf}{1.0}\PY{o}{+}\PY{l+m+mf}{0.0}\PY{n}{j}\PY{p}{)}\PY{p}{)}\PY{p}{)}\PY{p}{)}
\PY{c+c1}{\PYZsh{}print(\PYZdq{}(Table 2) Maximum error with ee interaction:   \PYZdq{}, np.max(np.abs(tbl002[\PYZdq{}New End\PYZdq{}].to\PYZus{}numpy()\PYZhy{}(1.0+0.0j))))}
\PY{c+c1}{\PYZsh{}print(\PYZdq{}(Table 2) Maximum error without ee interaction:\PYZdq{}, np.max(np.abs(tbl002[\PYZdq{}Exact End\PYZdq{}].to\PYZus{}numpy()\PYZhy{}(1.0+0.0j))))}
\PY{n+nb}{print}\PY{p}{(}\PY{l+s+s2}{\PYZdq{}}\PY{l+s+s2}{(Table 3) Maximum error with ee interaction:   }\PY{l+s+s2}{\PYZdq{}}\PY{p}{,} \PY{n}{np}\PY{o}{.}\PY{n}{max}\PY{p}{(}\PY{n}{np}\PY{o}{.}\PY{n}{abs}\PY{p}{(}\PY{n}{tbl003}\PY{p}{[}\PY{l+s+s2}{\PYZdq{}}\PY{l+s+s2}{New End}\PY{l+s+s2}{\PYZdq{}}\PY{p}{]}\PY{o}{.}\PY{n}{to\PYZus{}numpy}\PY{p}{(}\PY{p}{)}\PY{o}{\PYZhy{}}\PY{p}{(}\PY{l+m+mf}{1.0}\PY{o}{+}\PY{l+m+mf}{0.0}\PY{n}{j}\PY{p}{)}\PY{p}{)}\PY{p}{)}\PY{p}{)}
\PY{n+nb}{print}\PY{p}{(}\PY{l+s+s2}{\PYZdq{}}\PY{l+s+s2}{(Table 3) Maximum error without ee interaction:}\PY{l+s+s2}{\PYZdq{}}\PY{p}{,} \PY{n}{np}\PY{o}{.}\PY{n}{max}\PY{p}{(}\PY{n}{np}\PY{o}{.}\PY{n}{abs}\PY{p}{(}\PY{n}{tbl003}\PY{p}{[}\PY{l+s+s2}{\PYZdq{}}\PY{l+s+s2}{Exact End}\PY{l+s+s2}{\PYZdq{}}\PY{p}{]}\PY{o}{.}\PY{n}{to\PYZus{}numpy}\PY{p}{(}\PY{p}{)}\PY{o}{\PYZhy{}}\PY{p}{(}\PY{l+m+mf}{1.0}\PY{o}{+}\PY{l+m+mf}{0.0}\PY{n}{j}\PY{p}{)}\PY{p}{)}\PY{p}{)}\PY{p}{)}
\PY{c+c1}{\PYZsh{}print(\PYZdq{}(Table 4) Maximum error with ee interaction:   \PYZdq{}, np.max(np.abs(tbl004[\PYZdq{}New End\PYZdq{}].to\PYZus{}numpy()\PYZhy{}(1.0+0.0j))))}
\PY{c+c1}{\PYZsh{}print(\PYZdq{}(Table 4) Maximum error without ee interaction:\PYZdq{}, np.max(np.abs(tbl004[\PYZdq{}Exact End\PYZdq{}].to\PYZus{}numpy()\PYZhy{}(1.0+0.0j))))}
\PY{n+nb}{print}\PY{p}{(}\PY{l+s+s2}{\PYZdq{}}\PY{l+s+s2}{(Table 5) Maximum error with ee interaction:   }\PY{l+s+s2}{\PYZdq{}}\PY{p}{,} \PY{n}{np}\PY{o}{.}\PY{n}{max}\PY{p}{(}\PY{n}{np}\PY{o}{.}\PY{n}{abs}\PY{p}{(}\PY{n}{tbl005}\PY{p}{[}\PY{l+s+s2}{\PYZdq{}}\PY{l+s+s2}{New End}\PY{l+s+s2}{\PYZdq{}}\PY{p}{]}\PY{o}{.}\PY{n}{to\PYZus{}numpy}\PY{p}{(}\PY{p}{)}\PY{o}{\PYZhy{}}\PY{p}{(}\PY{l+m+mf}{1.0}\PY{o}{+}\PY{l+m+mf}{0.0}\PY{n}{j}\PY{p}{)}\PY{p}{)}\PY{p}{)}\PY{p}{)}
\PY{n+nb}{print}\PY{p}{(}\PY{l+s+s2}{\PYZdq{}}\PY{l+s+s2}{(Table 5) Maximum error without ee interaction:}\PY{l+s+s2}{\PYZdq{}}\PY{p}{,} \PY{n}{np}\PY{o}{.}\PY{n}{max}\PY{p}{(}\PY{n}{np}\PY{o}{.}\PY{n}{abs}\PY{p}{(}\PY{n}{tbl005}\PY{p}{[}\PY{l+s+s2}{\PYZdq{}}\PY{l+s+s2}{Exact End}\PY{l+s+s2}{\PYZdq{}}\PY{p}{]}\PY{o}{.}\PY{n}{to\PYZus{}numpy}\PY{p}{(}\PY{p}{)}\PY{o}{\PYZhy{}}\PY{p}{(}\PY{l+m+mf}{1.0}\PY{o}{+}\PY{l+m+mf}{0.0}\PY{n}{j}\PY{p}{)}\PY{p}{)}\PY{p}{)}\PY{p}{)}
\end{Verbatim}
\end{tcolorbox}

    \begin{Verbatim}[commandchars=\\\{\}]
(Table 0) Maximum error with ee interaction:    0.00998825640313818
(Table 0) Maximum error without ee interaction: 1.3073150373962005e-10
(Table 1) Maximum error with ee interaction:    0.00999628422114756
(Table 1) Maximum error without ee interaction: 2.470210660871276e-10
(Table 3) Maximum error with ee interaction:    0.009999118362851141
(Table 3) Maximum error without ee interaction: 6.86442258220229e-11
(Table 5) Maximum error with ee interaction:    0.00997567727459663
(Table 5) Maximum error without ee interaction: 1.7774866325471916e-09
    \end{Verbatim}

    Now to dump the unneeded data from the tables.

From previous analysis, I know that ``I v2'' is the best version of my
calculation for the current of the nonlinear solution. This value is
equal to the ``I v3'' current, which is calculated at our initial
position, using the full equation for \(I_j\), \begin{equation}
I_j = \frac{1}{2 k_0 i}\left( \overline{\psi_j} \frac{d}{dx} \psi_j -\psi_j \overline{\frac{d}{dx} \psi_j} -i \frac{2 B \pi R_j}{\Phi_0} \overline{\psi_j}  \psi_j  \right),
\end{equation} with our known \(\psi(0)\) and \(\psi^\prime(0)\) values.
We will thus remove ``I v1'' and ``I v3'' from our table.

\(k\) and \(A\) are unnecessary as they are always identical. After I
have checked ``Exact End'' and ``New End'', they are no longer required.
That leaves 4 independent variables and 11 dependent variables. Note
that here there is still unnecessary data, but at this point I don't
know it is unnecessary.

    \begin{tcolorbox}[breakable, size=fbox, boxrule=1pt, pad at break*=1mm,colback=cellbackground, colframe=cellborder]
\prompt{In}{incolor}{14}{\boxspacing}
\begin{Verbatim}[commandchars=\\\{\}]
\PY{n}{tbl2}   \PY{o}{=}    \PY{n}{tbl}\PY{p}{[}\PY{p}{[}\PY{l+s+s2}{\PYZdq{}}\PY{l+s+s2}{R}\PY{l+s+s2}{\PYZdq{}}\PY{p}{,} \PY{l+s+s2}{\PYZdq{}}\PY{l+s+s2}{B}\PY{l+s+s2}{\PYZdq{}}\PY{p}{,} \PY{l+s+s2}{\PYZdq{}}\PY{l+s+s2}{dk}\PY{l+s+s2}{\PYZdq{}}\PY{p}{,} \PY{l+s+s2}{\PYZdq{}}\PY{l+s+s2}{mu}\PY{l+s+s2}{\PYZdq{}}\PY{p}{,} \PY{l+s+s1}{\PYZsq{}}\PY{l+s+s1}{dpsi0}\PY{l+s+s1}{\PYZsq{}}\PY{p}{,} \PY{l+s+s1}{\PYZsq{}}\PY{l+s+s1}{a0}\PY{l+s+s1}{\PYZsq{}}\PY{p}{,} \PY{l+s+s1}{\PYZsq{}}\PY{l+s+s1}{b0}\PY{l+s+s1}{\PYZsq{}}\PY{p}{,} \PY{l+s+s1}{\PYZsq{}}\PY{l+s+s1}{A max 0}\PY{l+s+s1}{\PYZsq{}}\PY{p}{,} \PY{l+s+s1}{\PYZsq{}}\PY{l+s+s1}{I0}\PY{l+s+s1}{\PYZsq{}}\PY{p}{,} \PY{l+s+s1}{\PYZsq{}}\PY{l+s+s1}{dpsi}\PY{l+s+s1}{\PYZsq{}}\PY{p}{,} \PY{l+s+s1}{\PYZsq{}}\PY{l+s+s1}{a}\PY{l+s+s1}{\PYZsq{}}\PY{p}{,} \PY{l+s+s1}{\PYZsq{}}\PY{l+s+s1}{b}\PY{l+s+s1}{\PYZsq{}}\PY{p}{,} \PY{l+s+s1}{\PYZsq{}}\PY{l+s+s1}{A max new}\PY{l+s+s1}{\PYZsq{}}\PY{p}{,} \PY{l+s+s1}{\PYZsq{}}\PY{l+s+s1}{I v2}\PY{l+s+s1}{\PYZsq{}}\PY{p}{,} \PY{l+s+s1}{\PYZsq{}}\PY{l+s+s1}{effective mu}\PY{l+s+s1}{\PYZsq{}}\PY{p}{]}\PY{p}{]}
\PY{n}{tbl011} \PY{o}{=} \PY{n}{tbl001}\PY{p}{[}\PY{p}{[}\PY{l+s+s2}{\PYZdq{}}\PY{l+s+s2}{R}\PY{l+s+s2}{\PYZdq{}}\PY{p}{,} \PY{l+s+s2}{\PYZdq{}}\PY{l+s+s2}{B}\PY{l+s+s2}{\PYZdq{}}\PY{p}{,} \PY{l+s+s2}{\PYZdq{}}\PY{l+s+s2}{dk}\PY{l+s+s2}{\PYZdq{}}\PY{p}{,} \PY{l+s+s2}{\PYZdq{}}\PY{l+s+s2}{mu}\PY{l+s+s2}{\PYZdq{}}\PY{p}{,} \PY{l+s+s1}{\PYZsq{}}\PY{l+s+s1}{dpsi0}\PY{l+s+s1}{\PYZsq{}}\PY{p}{,} \PY{l+s+s1}{\PYZsq{}}\PY{l+s+s1}{a0}\PY{l+s+s1}{\PYZsq{}}\PY{p}{,} \PY{l+s+s1}{\PYZsq{}}\PY{l+s+s1}{b0}\PY{l+s+s1}{\PYZsq{}}\PY{p}{,} \PY{l+s+s1}{\PYZsq{}}\PY{l+s+s1}{A max 0}\PY{l+s+s1}{\PYZsq{}}\PY{p}{,} \PY{l+s+s1}{\PYZsq{}}\PY{l+s+s1}{I0}\PY{l+s+s1}{\PYZsq{}}\PY{p}{,} \PY{l+s+s1}{\PYZsq{}}\PY{l+s+s1}{dpsi}\PY{l+s+s1}{\PYZsq{}}\PY{p}{,} \PY{l+s+s1}{\PYZsq{}}\PY{l+s+s1}{a}\PY{l+s+s1}{\PYZsq{}}\PY{p}{,} \PY{l+s+s1}{\PYZsq{}}\PY{l+s+s1}{b}\PY{l+s+s1}{\PYZsq{}}\PY{p}{,} \PY{l+s+s1}{\PYZsq{}}\PY{l+s+s1}{A max new}\PY{l+s+s1}{\PYZsq{}}\PY{p}{,} \PY{l+s+s1}{\PYZsq{}}\PY{l+s+s1}{I v2}\PY{l+s+s1}{\PYZsq{}}\PY{p}{,} \PY{l+s+s1}{\PYZsq{}}\PY{l+s+s1}{effective mu}\PY{l+s+s1}{\PYZsq{}}\PY{p}{]}\PY{p}{]}
\PY{c+c1}{\PYZsh{}tbl012 = tbl002[[\PYZdq{}R\PYZdq{}, \PYZdq{}B\PYZdq{}, \PYZdq{}dk\PYZdq{}, \PYZdq{}mu\PYZdq{}, \PYZsq{}dpsi0\PYZsq{}, \PYZsq{}a0\PYZsq{}, \PYZsq{}b0\PYZsq{}, \PYZsq{}A max 0\PYZsq{}, \PYZsq{}I0\PYZsq{}, \PYZsq{}dpsi\PYZsq{}, \PYZsq{}a\PYZsq{}, \PYZsq{}b\PYZsq{}, \PYZsq{}A max new\PYZsq{}, \PYZsq{}I v2\PYZsq{}, \PYZsq{}effective mu\PYZsq{}]]}
\PY{n}{tbl013} \PY{o}{=} \PY{n}{tbl003}\PY{p}{[}\PY{p}{[}\PY{l+s+s2}{\PYZdq{}}\PY{l+s+s2}{R}\PY{l+s+s2}{\PYZdq{}}\PY{p}{,} \PY{l+s+s2}{\PYZdq{}}\PY{l+s+s2}{B}\PY{l+s+s2}{\PYZdq{}}\PY{p}{,} \PY{l+s+s2}{\PYZdq{}}\PY{l+s+s2}{dk}\PY{l+s+s2}{\PYZdq{}}\PY{p}{,} \PY{l+s+s2}{\PYZdq{}}\PY{l+s+s2}{mu}\PY{l+s+s2}{\PYZdq{}}\PY{p}{,} \PY{l+s+s1}{\PYZsq{}}\PY{l+s+s1}{dpsi0}\PY{l+s+s1}{\PYZsq{}}\PY{p}{,} \PY{l+s+s1}{\PYZsq{}}\PY{l+s+s1}{a0}\PY{l+s+s1}{\PYZsq{}}\PY{p}{,} \PY{l+s+s1}{\PYZsq{}}\PY{l+s+s1}{b0}\PY{l+s+s1}{\PYZsq{}}\PY{p}{,} \PY{l+s+s1}{\PYZsq{}}\PY{l+s+s1}{A max 0}\PY{l+s+s1}{\PYZsq{}}\PY{p}{,} \PY{l+s+s1}{\PYZsq{}}\PY{l+s+s1}{I0}\PY{l+s+s1}{\PYZsq{}}\PY{p}{,} \PY{l+s+s1}{\PYZsq{}}\PY{l+s+s1}{dpsi}\PY{l+s+s1}{\PYZsq{}}\PY{p}{,} \PY{l+s+s1}{\PYZsq{}}\PY{l+s+s1}{a}\PY{l+s+s1}{\PYZsq{}}\PY{p}{,} \PY{l+s+s1}{\PYZsq{}}\PY{l+s+s1}{b}\PY{l+s+s1}{\PYZsq{}}\PY{p}{,} \PY{l+s+s1}{\PYZsq{}}\PY{l+s+s1}{A max new}\PY{l+s+s1}{\PYZsq{}}\PY{p}{,} \PY{l+s+s1}{\PYZsq{}}\PY{l+s+s1}{I v2}\PY{l+s+s1}{\PYZsq{}}\PY{p}{,} \PY{l+s+s1}{\PYZsq{}}\PY{l+s+s1}{effective mu}\PY{l+s+s1}{\PYZsq{}}\PY{p}{]}\PY{p}{]}
\PY{c+c1}{\PYZsh{}tbl014 = tbl004[[\PYZdq{}R\PYZdq{}, \PYZdq{}B\PYZdq{}, \PYZdq{}dk\PYZdq{}, \PYZdq{}mu\PYZdq{}, \PYZsq{}dpsi0\PYZsq{}, \PYZsq{}a0\PYZsq{}, \PYZsq{}b0\PYZsq{}, \PYZsq{}A max 0\PYZsq{}, \PYZsq{}I0\PYZsq{}, \PYZsq{}dpsi\PYZsq{}, \PYZsq{}a\PYZsq{}, \PYZsq{}b\PYZsq{}, \PYZsq{}A max new\PYZsq{}, \PYZsq{}I v2\PYZsq{}, \PYZsq{}effective mu\PYZsq{}]]}
\PY{n}{tbl015} \PY{o}{=} \PY{n}{tbl005}\PY{p}{[}\PY{p}{[}\PY{l+s+s2}{\PYZdq{}}\PY{l+s+s2}{R}\PY{l+s+s2}{\PYZdq{}}\PY{p}{,} \PY{l+s+s2}{\PYZdq{}}\PY{l+s+s2}{B}\PY{l+s+s2}{\PYZdq{}}\PY{p}{,} \PY{l+s+s2}{\PYZdq{}}\PY{l+s+s2}{dk}\PY{l+s+s2}{\PYZdq{}}\PY{p}{,} \PY{l+s+s2}{\PYZdq{}}\PY{l+s+s2}{mu}\PY{l+s+s2}{\PYZdq{}}\PY{p}{,} \PY{l+s+s1}{\PYZsq{}}\PY{l+s+s1}{dpsi0}\PY{l+s+s1}{\PYZsq{}}\PY{p}{,} \PY{l+s+s1}{\PYZsq{}}\PY{l+s+s1}{a0}\PY{l+s+s1}{\PYZsq{}}\PY{p}{,} \PY{l+s+s1}{\PYZsq{}}\PY{l+s+s1}{b0}\PY{l+s+s1}{\PYZsq{}}\PY{p}{,} \PY{l+s+s1}{\PYZsq{}}\PY{l+s+s1}{A max 0}\PY{l+s+s1}{\PYZsq{}}\PY{p}{,} \PY{l+s+s1}{\PYZsq{}}\PY{l+s+s1}{I0}\PY{l+s+s1}{\PYZsq{}}\PY{p}{,} \PY{l+s+s1}{\PYZsq{}}\PY{l+s+s1}{dpsi}\PY{l+s+s1}{\PYZsq{}}\PY{p}{,} \PY{l+s+s1}{\PYZsq{}}\PY{l+s+s1}{a}\PY{l+s+s1}{\PYZsq{}}\PY{p}{,} \PY{l+s+s1}{\PYZsq{}}\PY{l+s+s1}{b}\PY{l+s+s1}{\PYZsq{}}\PY{p}{,} \PY{l+s+s1}{\PYZsq{}}\PY{l+s+s1}{A max new}\PY{l+s+s1}{\PYZsq{}}\PY{p}{,} \PY{l+s+s1}{\PYZsq{}}\PY{l+s+s1}{I v2}\PY{l+s+s1}{\PYZsq{}}\PY{p}{,} \PY{l+s+s1}{\PYZsq{}}\PY{l+s+s1}{effective mu}\PY{l+s+s1}{\PYZsq{}}\PY{p}{]}\PY{p}{]}
\end{Verbatim}
\end{tcolorbox}

    The next step is to clear the table of all the unneeded high energy
solutions, as I need to only keep the lowest energy solution for every
set of parameters. The functions to clear out this data are in the
table\_scripts.py file in my library.

    \begin{tcolorbox}[breakable, size=fbox, boxrule=1pt, pad at break*=1mm,colback=cellbackground, colframe=cellborder]
\prompt{In}{incolor}{17}{\boxspacing}
\begin{Verbatim}[commandchars=\\\{\}]
\PY{n}{tbl021} \PY{o}{=} \PY{n}{clean\PYZus{}table}\PY{p}{(}\PY{n}{tbl011}\PY{p}{)}
\PY{n}{tbl023} \PY{o}{=} \PY{n}{clean\PYZus{}table}\PY{p}{(}\PY{n}{tbl013}\PY{p}{)}
\end{Verbatim}
\end{tcolorbox}

    \begin{tcolorbox}[breakable, size=fbox, boxrule=1pt, pad at break*=1mm,colback=cellbackground, colframe=cellborder]
\prompt{In}{incolor}{18}{\boxspacing}
\begin{Verbatim}[commandchars=\\\{\}]
\PY{c+c1}{\PYZsh{} Monte Carlo data must be done with the Monte Carlo version of the algorithm, due to time constraints.}

\PY{c+c1}{\PYZsh{}tbl022 = clean\PYZus{}table\PYZus{}MC(tbl012)}
\PY{c+c1}{\PYZsh{}tbl024 = clean\PYZus{}table\PYZus{}MC(tbl014)}
\PY{n}{tbl025} \PY{o}{=} \PY{n}{clean\PYZus{}table\PYZus{}MC}\PY{p}{(}\PY{n}{tbl015}\PY{p}{)}
\end{Verbatim}
\end{tcolorbox}

    \begin{tcolorbox}[breakable, size=fbox, boxrule=1pt, pad at break*=1mm,colback=cellbackground, colframe=cellborder]
\prompt{In}{incolor}{19}{\boxspacing}
\begin{Verbatim}[commandchars=\\\{\}]
\PY{n}{new\PYZus{}tbl} \PY{o}{=} \PY{n}{clean\PYZus{}table}\PY{p}{(}\PY{n}{tbl2}\PY{p}{)}
\end{Verbatim}
\end{tcolorbox}

    And a test of just one of these dataframes.

    \begin{tcolorbox}[breakable, size=fbox, boxrule=1pt, pad at break*=1mm,colback=cellbackground, colframe=cellborder]
\prompt{In}{incolor}{20}{\boxspacing}
\begin{Verbatim}[commandchars=\\\{\}]
\PY{n}{tbl021}\PY{o}{.}\PY{n}{head}\PY{p}{(}\PY{p}{)}
\end{Verbatim}
\end{tcolorbox}

            \begin{tcolorbox}[breakable, size=fbox, boxrule=.5pt, pad at break*=1mm, opacityfill=0]
\prompt{Out}{outcolor}{20}{\boxspacing}
\begin{Verbatim}[commandchars=\\\{\}]
     R    B    dk            mu  \textbackslash{}
0  0.9  0.5  0.51  1.000000e-10
1  0.9  0.5  0.51  1.584893e-10
2  0.9  0.5  0.51  2.511886e-10
3  0.9  0.5  0.51  3.981072e-10
4  0.9  0.5  0.51  6.309573e-10

                                            dpsi0                  a0  \textbackslash{}
0  1.054613e+10+4.462116e+                    07j  0.501674-0.439413j
1  1.054613e+10+4.462116e+                    07j  0.501674-0.439413j
2  1.054613e+10+4.462116e+                    07j  0.501674-0.439413j
3  1.054613e+10+4.462116e+                    07j  0.501674-0.439413j
4  1.054613e+10+4.462116e+                    07j  0.501674-0.439413j

                   b0   A max 0        I0  \textbackslash{}
0  0.498326+0.439413j  1.331292  0.003348
1  0.498326+0.439413j  1.331292  0.003348
2  0.498326+0.439413j  1.331292  0.003348
3  0.498326+0.439413j  1.331292  0.003348
4  0.498326+0.439413j  1.331292  0.003348

                                             dpsi                   a  \textbackslash{}
0  1.052821e+10+4.460161e+                    07j  0.501673-0.438666j
1  1.051777e+10+4.459022e+                    07j  0.501673-0.438231j
2  1.050128e+10+4.457225e+                    07j  0.501672-0.437544j
3  1.047527e+10+4.454395e+                    07j  0.501671-0.436460j
4  1.043440e+10+4.449953e+                    07j  0.501669-0.434757j

                    b  A max new      I v2  effective mu
0  0.498327+0.438666j   1.330307  0.003346  1.769717e-10
1  0.498327+0.438231j   1.329733  0.003345  2.802394e-10
2  0.498328+0.437544j   1.328828  0.003344  4.435448e-10
3  0.498329+0.436460j   1.327402  0.003342  7.014631e-10
4  0.498331+0.434757j   1.325164  0.003338  1.107999e-09
\end{Verbatim}
\end{tcolorbox}
        
    \hypertarget{triangle-plots}{%
\section{Triangle Plots}\label{triangle-plots}}

    Now for the triangle plots. This is section is designed to get a sense
of what we are dealing with.

    \begin{tcolorbox}[breakable, size=fbox, boxrule=1pt, pad at break*=1mm,colback=cellbackground, colframe=cellborder]
\prompt{In}{incolor}{21}{\boxspacing}
\begin{Verbatim}[commandchars=\\\{\}]
\PY{n}{g} \PY{o}{=} \PY{n}{sns}\PY{o}{.}\PY{n}{PairGrid}\PY{p}{(}\PY{n}{new\PYZus{}tbl}\PY{p}{[}\PY{p}{[}\PY{l+s+s2}{\PYZdq{}}\PY{l+s+s2}{R}\PY{l+s+s2}{\PYZdq{}}\PY{p}{,} \PY{l+s+s2}{\PYZdq{}}\PY{l+s+s2}{B}\PY{l+s+s2}{\PYZdq{}}\PY{p}{,} \PY{l+s+s2}{\PYZdq{}}\PY{l+s+s2}{dk}\PY{l+s+s2}{\PYZdq{}}\PY{p}{,} \PY{l+s+s2}{\PYZdq{}}\PY{l+s+s2}{mu}\PY{l+s+s2}{\PYZdq{}}\PY{p}{,} \PY{l+s+s2}{\PYZdq{}}\PY{l+s+s2}{I0}\PY{l+s+s2}{\PYZdq{}}\PY{p}{,} \PY{l+s+s2}{\PYZdq{}}\PY{l+s+s2}{I v2}\PY{l+s+s2}{\PYZdq{}}\PY{p}{]}\PY{p}{]}\PY{p}{)}

\PY{n}{g}\PY{o}{.}\PY{n}{map\PYZus{}diag}\PY{p}{(}\PY{n}{sns}\PY{o}{.}\PY{n}{kdeplot}\PY{p}{)}
\PY{n}{g}\PY{o}{.}\PY{n}{map\PYZus{}upper}\PY{p}{(}\PY{n}{sns}\PY{o}{.}\PY{n}{scatterplot}\PY{p}{)}
\PY{n}{g}\PY{o}{.}\PY{n}{map\PYZus{}lower}\PY{p}{(}\PY{n}{sns}\PY{o}{.}\PY{n}{scatterplot}\PY{p}{)}
\end{Verbatim}
\end{tcolorbox}

            \begin{tcolorbox}[breakable, size=fbox, boxrule=.5pt, pad at break*=1mm, opacityfill=0]
\prompt{Out}{outcolor}{21}{\boxspacing}
\begin{Verbatim}[commandchars=\\\{\}]
<seaborn.axisgrid.PairGrid at 0x27b8be76c00>
\end{Verbatim}
\end{tcolorbox}
        
    \begin{center}
    \adjustimage{max size={0.9\linewidth}{0.9\paperheight}}{figs/bvp_fin_en/output_39_1.png}
    \end{center}
    { \hspace*{\fill} \\}
    
    \(\mu\) dependence reduces the current all-around for larger \(\mu\)
values. dk changes here are too sensitive, and the change is not visible
at all at this level of change. \(B\) changes are approximated by a
first order Taylor expansion -- the changes here are noticeable, but
linearized. The first derivative is not small, but we are still in the
linearized regime.

Potential paths for analysis: 1) Redo the grid for points with better
spacing. 2) Freeze \(R\) and analyze dependence on only the other
parameters. My understanding of the old current dependence is that \(R\)
dependence is more complex, but coupled with other dependencies. \(\mu\)
appears to cause a rotation of my functions as well as a reduction of
amplitude. If the coupling is unchanged, in theory determining \(B\), dk
and \(\mu\) dependence is enough for deciphering the \(R\) dependence 3)
Monte Carlo gathering of points and a regression to a postulated
functional form.

The \(I0\) vs ``\(I\) v2'' graph looks like there is a rotation.
Specifically, \(R\) is causing the remaining pattern to rotate and
rescale. The fact that it is specifically \(R\) which is causing the
wildly differing data spread will be seen later.

    \begin{tcolorbox}[breakable, size=fbox, boxrule=1pt, pad at break*=1mm,colback=cellbackground, colframe=cellborder]
\prompt{In}{incolor}{22}{\boxspacing}
\begin{Verbatim}[commandchars=\\\{\}]
\PY{n}{g} \PY{o}{=} \PY{n}{sns}\PY{o}{.}\PY{n}{PairGrid}\PY{p}{(}\PY{n}{tbl021}\PY{p}{[}\PY{p}{[}\PY{l+s+s2}{\PYZdq{}}\PY{l+s+s2}{R}\PY{l+s+s2}{\PYZdq{}}\PY{p}{,} \PY{l+s+s2}{\PYZdq{}}\PY{l+s+s2}{B}\PY{l+s+s2}{\PYZdq{}}\PY{p}{,} \PY{l+s+s2}{\PYZdq{}}\PY{l+s+s2}{dk}\PY{l+s+s2}{\PYZdq{}}\PY{p}{,} \PY{l+s+s2}{\PYZdq{}}\PY{l+s+s2}{mu}\PY{l+s+s2}{\PYZdq{}}\PY{p}{,} \PY{l+s+s2}{\PYZdq{}}\PY{l+s+s2}{I0}\PY{l+s+s2}{\PYZdq{}}\PY{p}{,} \PY{l+s+s2}{\PYZdq{}}\PY{l+s+s2}{I v2}\PY{l+s+s2}{\PYZdq{}}\PY{p}{]}\PY{p}{]}\PY{p}{)}

\PY{n}{g}\PY{o}{.}\PY{n}{map\PYZus{}diag}\PY{p}{(}\PY{n}{sns}\PY{o}{.}\PY{n}{kdeplot}\PY{p}{)}
\PY{n}{g}\PY{o}{.}\PY{n}{map\PYZus{}upper}\PY{p}{(}\PY{n}{sns}\PY{o}{.}\PY{n}{scatterplot}\PY{p}{)}
\PY{n}{g}\PY{o}{.}\PY{n}{map\PYZus{}lower}\PY{p}{(}\PY{n}{sns}\PY{o}{.}\PY{n}{scatterplot}\PY{p}{)}
\end{Verbatim}
\end{tcolorbox}

            \begin{tcolorbox}[breakable, size=fbox, boxrule=.5pt, pad at break*=1mm, opacityfill=0]
\prompt{Out}{outcolor}{22}{\boxspacing}
\begin{Verbatim}[commandchars=\\\{\}]
<seaborn.axisgrid.PairGrid at 0x27b983ae420>
\end{Verbatim}
\end{tcolorbox}
        
    \begin{center}
    \adjustimage{max size={0.9\linewidth}{0.9\paperheight}}{figs/bvp_fin_en/output_41_1.png}
    \end{center}
    { \hspace*{\fill} \\}
    
    This is the same data set as the previous graph, just now with more
points. I see a dk vs ``I v2'' jump, which I didn't see previously.

    \begin{tcolorbox}[breakable, size=fbox, boxrule=1pt, pad at break*=1mm,colback=cellbackground, colframe=cellborder]
\prompt{In}{incolor}{23}{\boxspacing}
\begin{Verbatim}[commandchars=\\\{\}]
\PY{n}{g} \PY{o}{=} \PY{n}{sns}\PY{o}{.}\PY{n}{PairGrid}\PY{p}{(}\PY{n}{tbl023}\PY{p}{[}\PY{p}{[}\PY{l+s+s2}{\PYZdq{}}\PY{l+s+s2}{R}\PY{l+s+s2}{\PYZdq{}}\PY{p}{,} \PY{l+s+s2}{\PYZdq{}}\PY{l+s+s2}{B}\PY{l+s+s2}{\PYZdq{}}\PY{p}{,} \PY{l+s+s2}{\PYZdq{}}\PY{l+s+s2}{dk}\PY{l+s+s2}{\PYZdq{}}\PY{p}{,} \PY{l+s+s2}{\PYZdq{}}\PY{l+s+s2}{mu}\PY{l+s+s2}{\PYZdq{}}\PY{p}{,} \PY{l+s+s2}{\PYZdq{}}\PY{l+s+s2}{I0}\PY{l+s+s2}{\PYZdq{}}\PY{p}{,} \PY{l+s+s2}{\PYZdq{}}\PY{l+s+s2}{I v2}\PY{l+s+s2}{\PYZdq{}}\PY{p}{]}\PY{p}{]}\PY{p}{)}

\PY{n}{g}\PY{o}{.}\PY{n}{map\PYZus{}diag}\PY{p}{(}\PY{n}{sns}\PY{o}{.}\PY{n}{kdeplot}\PY{p}{)}
\PY{n}{g}\PY{o}{.}\PY{n}{map\PYZus{}upper}\PY{p}{(}\PY{n}{sns}\PY{o}{.}\PY{n}{scatterplot}\PY{p}{)}
\PY{n}{g}\PY{o}{.}\PY{n}{map\PYZus{}lower}\PY{p}{(}\PY{n}{sns}\PY{o}{.}\PY{n}{scatterplot}\PY{p}{)}
\end{Verbatim}
\end{tcolorbox}

            \begin{tcolorbox}[breakable, size=fbox, boxrule=.5pt, pad at break*=1mm, opacityfill=0]
\prompt{Out}{outcolor}{23}{\boxspacing}
\begin{Verbatim}[commandchars=\\\{\}]
<seaborn.axisgrid.PairGrid at 0x27b9935e420>
\end{Verbatim}
\end{tcolorbox}
        
    \begin{center}
    \adjustimage{max size={0.9\linewidth}{0.9\paperheight}}{figs/bvp_fin_en/output_43_1.png}
    \end{center}
    { \hspace*{\fill} \\}
    
    As expected, analyzing the results for large changes in \(R\) and \(B\)
just causes confusion. \(\mu\) dependence and \(dk\) dependence are
quite clear here. \(\mu\) appears to be reducing the amplitude of the
current.

    \begin{tcolorbox}[breakable, size=fbox, boxrule=1pt, pad at break*=1mm,colback=cellbackground, colframe=cellborder]
\prompt{In}{incolor}{51}{\boxspacing}
\begin{Verbatim}[commandchars=\\\{\}]
\PY{n}{g} \PY{o}{=} \PY{n}{sns}\PY{o}{.}\PY{n}{PairGrid}\PY{p}{(}\PY{n}{tbl025}\PY{p}{[}\PY{p}{[}\PY{l+s+s2}{\PYZdq{}}\PY{l+s+s2}{R}\PY{l+s+s2}{\PYZdq{}}\PY{p}{,} \PY{l+s+s2}{\PYZdq{}}\PY{l+s+s2}{B}\PY{l+s+s2}{\PYZdq{}}\PY{p}{,} \PY{l+s+s2}{\PYZdq{}}\PY{l+s+s2}{dk}\PY{l+s+s2}{\PYZdq{}}\PY{p}{,} \PY{l+s+s2}{\PYZdq{}}\PY{l+s+s2}{mu}\PY{l+s+s2}{\PYZdq{}}\PY{p}{,} \PY{l+s+s2}{\PYZdq{}}\PY{l+s+s2}{I0}\PY{l+s+s2}{\PYZdq{}}\PY{p}{,} \PY{l+s+s2}{\PYZdq{}}\PY{l+s+s2}{I v2}\PY{l+s+s2}{\PYZdq{}}\PY{p}{]}\PY{p}{]}\PY{p}{)}

\PY{n}{g}\PY{o}{.}\PY{n}{map\PYZus{}diag}\PY{p}{(}\PY{n}{sns}\PY{o}{.}\PY{n}{kdeplot}\PY{p}{)}
\PY{n}{g}\PY{o}{.}\PY{n}{map\PYZus{}upper}\PY{p}{(}\PY{n}{sns}\PY{o}{.}\PY{n}{scatterplot}\PY{p}{)}
\PY{n}{g}\PY{o}{.}\PY{n}{map\PYZus{}lower}\PY{p}{(}\PY{n}{sns}\PY{o}{.}\PY{n}{scatterplot}\PY{p}{)}
\end{Verbatim}
\end{tcolorbox}

            \begin{tcolorbox}[breakable, size=fbox, boxrule=.5pt, pad at break*=1mm, opacityfill=0]
\prompt{Out}{outcolor}{51}{\boxspacing}
\begin{Verbatim}[commandchars=\\\{\}]
<seaborn.axisgrid.PairGrid at 0x19093d11f70>
\end{Verbatim}
\end{tcolorbox}
        
    \begin{center}
    \adjustimage{max size={0.9\linewidth}{0.9\paperheight}}{figs/bvp_fin_en/output_45_1.png}
    \end{center}
    { \hspace*{\fill} \\}
    
    The Monte Carlo analysis has an error. \(\mu\) is fixed here at
\(10^{-7}\). This was not intended. Certain shapes and patterns are
visible here. \(B\) dependence is linear (or rather linearized). \(R\)
dependence of I0 shows a divergence. My choices of \(B\) and \(R\)
appear to be significant with respect to the part of the periodic
behavior they show.

    \hypertarget{fixed-r-analysis}{%
\section{\texorpdfstring{Fixed \(R\)
Analysis}{Fixed R Analysis}}\label{fixed-r-analysis}}

    My first attempt at analysis fixed \(R\) and attempted to analyze the
other parameters from the grid with small parameter variations. \(R\) is
the most difficult of the parameters to analyze, so I removed it to try
to decipher the other dependencies. Note that I expected the form to be
similar, but not identical, to the case without ee interaction. (It
turns out to be identical.) Here I retain the graphs, but none of the
analysis (as it is obsolete).

    \begin{tcolorbox}[breakable, size=fbox, boxrule=1pt, pad at break*=1mm,colback=cellbackground, colframe=cellborder]
\prompt{In}{incolor}{24}{\boxspacing}
\begin{Verbatim}[commandchars=\\\{\}]
\PY{n}{R\PYZus{}val} \PY{o}{=} \PY{l+m+mf}{0.900018} \PY{o}{*} \PY{l+m+mf}{1.0e\PYZhy{}6}
\PY{n}{tbl031} \PY{o}{=} \PY{n}{tbl021}\PY{p}{[}\PY{n}{tbl021}\PY{p}{[}\PY{l+s+s2}{\PYZdq{}}\PY{l+s+s2}{R}\PY{l+s+s2}{\PYZdq{}}\PY{p}{]}\PY{o}{==} \PY{l+m+mf}{0.900018}\PY{p}{]}
\end{Verbatim}
\end{tcolorbox}

    \begin{tcolorbox}[breakable, size=fbox, boxrule=1pt, pad at break*=1mm,colback=cellbackground, colframe=cellborder]
\prompt{In}{incolor}{25}{\boxspacing}
\begin{Verbatim}[commandchars=\\\{\}]
\PY{n}{g} \PY{o}{=} \PY{n}{sns}\PY{o}{.}\PY{n}{PairGrid}\PY{p}{(}\PY{n}{tbl031}\PY{p}{[}\PY{p}{[}\PY{l+s+s2}{\PYZdq{}}\PY{l+s+s2}{B}\PY{l+s+s2}{\PYZdq{}}\PY{p}{,} \PY{l+s+s2}{\PYZdq{}}\PY{l+s+s2}{dk}\PY{l+s+s2}{\PYZdq{}}\PY{p}{,} \PY{l+s+s2}{\PYZdq{}}\PY{l+s+s2}{mu}\PY{l+s+s2}{\PYZdq{}}\PY{p}{,} \PY{l+s+s2}{\PYZdq{}}\PY{l+s+s2}{I0}\PY{l+s+s2}{\PYZdq{}}\PY{p}{,} \PY{l+s+s2}{\PYZdq{}}\PY{l+s+s2}{I v2}\PY{l+s+s2}{\PYZdq{}}\PY{p}{]}\PY{p}{]}\PY{p}{)}

\PY{n}{g}\PY{o}{.}\PY{n}{map\PYZus{}diag}\PY{p}{(}\PY{n}{sns}\PY{o}{.}\PY{n}{kdeplot}\PY{p}{)}
\PY{n}{g}\PY{o}{.}\PY{n}{map\PYZus{}upper}\PY{p}{(}\PY{n}{sns}\PY{o}{.}\PY{n}{scatterplot}\PY{p}{)}
\PY{n}{g}\PY{o}{.}\PY{n}{map\PYZus{}lower}\PY{p}{(}\PY{n}{sns}\PY{o}{.}\PY{n}{scatterplot}\PY{p}{)}
\end{Verbatim}
\end{tcolorbox}

            \begin{tcolorbox}[breakable, size=fbox, boxrule=.5pt, pad at break*=1mm, opacityfill=0]
\prompt{Out}{outcolor}{25}{\boxspacing}
\begin{Verbatim}[commandchars=\\\{\}]
<seaborn.axisgrid.PairGrid at 0x27b9593dd00>
\end{Verbatim}
\end{tcolorbox}
        
    \begin{center}
    \adjustimage{max size={0.9\linewidth}{0.9\paperheight}}{figs/bvp_fin_en/output_50_1.png}
    \end{center}
    { \hspace*{\fill} \\}
    
    My idea for the model of \(I_j\) is to use the old \(I_j\) calculations,
defining a phase shift given by the change in \(k\) and \(M\), and then
adding a power-law squeezing of the current based on \(\mu\). The
squashing and rotation are clearly visible on the plots, with the
squashing being visible even with the rapid rotation. I want to start
with the predicted changes of \(k\) and \(M\) to guess my phase shift as
everything appears to be based on this change, and I know that they will
produce a phase shift. I should probably look not just power-law decay
for \(\mu\). I am not quite sure why I just jumped to power law.
Possibly because power laws are typical of emergent behavior.

    \hypertarget{predicting-functional-form}{%
\section{Predicting Functional Form}\label{predicting-functional-form}}

    Now to add our model to our table and see how well it predicts our
results. Our model is, \begin{equation}
I_j = - \frac{\sin\left(2 \pi M_j R_j\right)}{\sin\left(2 \pi k_j R_j\right)}
\end{equation}

    \begin{tcolorbox}[breakable, size=fbox, boxrule=1pt, pad at break*=1mm,colback=cellbackground, colframe=cellborder]
\prompt{In}{incolor}{26}{\boxspacing}
\begin{Verbatim}[commandchars=\\\{\}]
\PY{c+c1}{\PYZsh{} Lets predict a model which looks like the model without ee interaction, and see if this provides}
\PY{c+c1}{\PYZsh{} an easier starting point for regression.}
\PY{c+c1}{\PYZsh{} Note that R and B are here both a percent of their maximum value, not in SI units.}

\PY{k}{def} \PY{n+nf}{func\PYZus{}cur\PYZus{}ee}\PY{p}{(}\PY{n}{R}\PY{p}{,} \PY{n}{B}\PY{p}{,} \PY{n}{dk}\PY{p}{,} \PY{n}{mu}\PY{p}{,} \PY{n}{dpsi0}\PY{p}{,} \PY{n}{A} \PY{o}{=} \PY{l+m+mf}{1.}\PY{p}{,} \PY{n}{k} \PY{o}{=} \PY{n}{eelib}\PY{o}{.}\PY{n}{kFAu}\PY{p}{)}\PY{p}{:}
    \PY{n}{R} \PY{o}{=} \PY{n}{R} \PY{o}{*} \PY{n}{R\PYZus{}max}
    \PY{n}{B} \PY{o}{=} \PY{n}{B} \PY{o}{*} \PY{n}{B\PYZus{}max}
    \PY{n}{k\PYZus{}new} \PY{o}{=} \PY{n}{eelib}\PY{o}{.}\PY{n}{pred\PYZus{}fast\PYZus{}k\PYZus{}true}\PY{p}{(}\PY{n}{dpsi0}\PY{p}{,} \PY{n}{mu}\PY{p}{,} \PY{n}{dk}\PY{p}{,} \PY{n}{B}\PY{p}{,} \PY{n}{R}\PY{p}{,} \PY{n}{A}\PY{o}{=}\PY{n}{A}\PY{p}{,} \PY{n}{k0}\PY{o}{=}\PY{n}{k}\PY{p}{)}
    \PY{n}{M\PYZus{}new} \PY{o}{=} \PY{n}{eelib}\PY{o}{.}\PY{n}{pred\PYZus{}slow\PYZus{}k\PYZus{}v3}\PY{p}{(}\PY{n}{dpsi0}\PY{p}{,} \PY{n}{mu}\PY{p}{,} \PY{n}{dk}\PY{p}{,} \PY{n}{B}\PY{p}{,} \PY{n}{R}\PY{p}{,} \PY{n}{A}\PY{o}{=}\PY{n}{A}\PY{p}{,} \PY{n}{k0}\PY{o}{=}\PY{n}{k}\PY{p}{)}
    
    \PY{n}{ret\PYZus{}val} \PY{o}{=} \PY{o}{\PYZhy{}} \PY{n}{np}\PY{o}{.}\PY{n}{sin}\PY{p}{(}\PY{l+m+mi}{2}\PY{o}{*}\PY{n}{pi}\PY{o}{*}\PY{n}{R}\PY{o}{*}\PY{n}{M\PYZus{}new}\PY{p}{)}\PY{o}{/}\PY{n}{np}\PY{o}{.}\PY{n}{sin}\PY{p}{(}\PY{l+m+mi}{2}\PY{o}{*}\PY{n}{pi}\PY{o}{*}\PY{n}{R}\PY{o}{*}\PY{n}{k\PYZus{}new}\PY{p}{)}

    \PY{k}{return} \PY{n}{ret\PYZus{}val}
\end{Verbatim}
\end{tcolorbox}

    And for the linear case, \begin{equation}
I_j = - \frac{\sin\left(2 \pi M_0 R_j\right)}{\sin\left(2 \pi k_0 R_j\right)}
\end{equation}

    \begin{tcolorbox}[breakable, size=fbox, boxrule=1pt, pad at break*=1mm,colback=cellbackground, colframe=cellborder]
\prompt{In}{incolor}{ }{\boxspacing}
\begin{Verbatim}[commandchars=\\\{\}]
\PY{c+c1}{\PYZsh{} Our known solution for the linear case. }
\PY{c+c1}{\PYZsh{} Note that R and B are here both a percent of their maximum value, not in SI units.}

\PY{k}{def} \PY{n+nf}{func\PYZus{}cur}\PY{p}{(}\PY{n}{R}\PY{p}{,} \PY{n}{B}\PY{p}{,} \PY{n}{k}\PY{p}{,} \PY{n}{dk}\PY{p}{)}\PY{p}{:}
    \PY{n}{R} \PY{o}{=} \PY{n}{R} \PY{o}{*} \PY{n}{R\PYZus{}max}
    \PY{n}{B} \PY{o}{=} \PY{n}{B} \PY{o}{*} \PY{n}{B\PYZus{}max}
    \PY{n}{ret\PYZus{}val} \PY{o}{=} \PY{o}{\PYZhy{}} \PY{n}{np}\PY{o}{.}\PY{n}{sin}\PY{p}{(}\PY{l+m+mi}{2}\PY{o}{*}\PY{n}{pi}\PY{o}{*}\PY{n}{R}\PY{o}{*}\PY{n}{R}\PY{o}{*}\PY{n}{B}\PY{o}{*}\PY{n}{phi0inv}\PY{p}{)}\PY{o}{/}\PY{n}{np}\PY{o}{.}\PY{n}{sin}\PY{p}{(}\PY{l+m+mi}{2}\PY{o}{*}\PY{n}{pi}\PY{o}{*}\PY{n}{R}\PY{o}{*}\PY{p}{(}\PY{n}{k}\PY{o}{+}\PY{n}{dk}\PY{o}{/}\PY{l+m+mi}{2}\PY{o}{/}\PY{n}{R\PYZus{}max}\PY{p}{)}\PY{p}{)}
    \PY{k}{return} \PY{n}{ret\PYZus{}val}
\end{Verbatim}
\end{tcolorbox}

    I had an error in my code, so in order to try to fix it, I tried to
develop an alternative equation for calculating current, in order to try
to reduce numerical error. It is derived as follows. I know that
\(a_j+b_j = \psi(0)\) which is set to \(1\). I also know that
\(I_j = |a_j|^2-|b_j|^2\). So, \begin{equation}
a_j-b_j = \bar{\psi}_j(0) * (a_j-b_j) =(\bar{a}_j+\bar{b}_j)(a_j-b_j) =|a_j|^2-|b_j|^2 +\bar{a}_jb_j-a_j\bar{b}_j = I - \bar{a}_j(1-a_j)+a(1-\bar{a}_j) = I + 2 i \, \mathrm{Im} (a_j).
\end{equation} We can also derive an equation for \(a_j-b_j\) from the
derivative of \(\psi_j\). \begin{equation}
\psi_j^\prime(x) = i a_j (k_j+M_j) e^{i (k_j+M_j) x} + i b_j (M_j-_jk) e^{i (M_j-k_j) x} 
\end{equation} or for \(x=0\), \begin{equation}
\psi_j^\prime(0) = i M_j + i k_j \left( a_j - b_j \right).
\end{equation} So, \begin{equation}
a_j-b_j = 2a_j-1 = \frac{\psi_j^\prime(0) - i M_j}{ik_j}
\end{equation} and \begin{equation}
a_j = - i \frac{\psi_j^\prime(0)}{2 k_j}-\frac{M_j}{2 k_j}+\frac{1}{2}.
\end{equation} Putting everything together gives \begin{equation}
I_j = \frac{\mathrm{Im} \left(\psi_j^\prime(0)\right)}{k_j} - \frac{M_j}{k_j}
\end{equation}

    \begin{tcolorbox}[breakable, size=fbox, boxrule=1pt, pad at break*=1mm,colback=cellbackground, colframe=cellborder]
\prompt{In}{incolor}{ }{\boxspacing}
\begin{Verbatim}[commandchars=\\\{\}]
\PY{c+c1}{\PYZsh{} New current calculation.}
\PY{k}{def} \PY{n+nf}{func\PYZus{}cur\PYZus{}from\PYZus{}deriv}\PY{p}{(}\PY{n}{dpsi0}\PY{p}{,} \PY{n}{M}\PY{p}{,} \PY{n}{k}\PY{p}{)}\PY{p}{:}
    \PY{k}{return} \PY{n}{np}\PY{o}{.}\PY{n}{imag}\PY{p}{(}\PY{n}{dpsi0}\PY{p}{)}\PY{o}{/}\PY{n}{k}\PY{o}{\PYZhy{}}\PY{n}{M}\PY{o}{/}\PY{n}{k}
\end{Verbatim}
\end{tcolorbox}

    Now to add these to the table.

    \begin{tcolorbox}[breakable, size=fbox, boxrule=1pt, pad at break*=1mm,colback=cellbackground, colframe=cellborder]
\prompt{In}{incolor}{36}{\boxspacing}
\begin{Verbatim}[commandchars=\\\{\}]
\PY{c+c1}{\PYZsh{} Code for adding columns to a table with the predictions for}
\PY{c+c1}{\PYZsh{} I0 and I. }
\PY{k}{def} \PY{n+nf}{predict\PYZus{}table}\PY{p}{(}\PY{n}{table}\PY{p}{,} \PY{n}{k\PYZus{}val}\PY{o}{=}\PY{n}{eelib}\PY{o}{.}\PY{n}{kFAu}\PY{p}{,} \PY{n}{A\PYZus{}val}\PY{o}{=}\PY{l+m+mf}{1.0}\PY{p}{)}\PY{p}{:}
    \PY{c+c1}{\PYZsh{} Find out how many rows we have.}
    \PY{n}{no\PYZus{}samp} \PY{o}{=} \PY{n+nb}{len}\PY{p}{(}\PY{n}{table}\PY{o}{.}\PY{n}{index}\PY{p}{)}

    \PY{c+c1}{\PYZsh{} Storage arrays, preallocate our space}
    \PY{n}{pred\PYZus{}vals} \PY{o}{=} \PY{n}{np}\PY{o}{.}\PY{n}{zeros}\PY{p}{(}\PY{n}{no\PYZus{}samp}\PY{p}{)} \PY{c+c1}{\PYZsh{} prediction for I0 }
    \PY{n}{cur\PYZus{}vals} \PY{o}{=} \PY{n}{np}\PY{o}{.}\PY{n}{zeros}\PY{p}{(}\PY{n}{no\PYZus{}samp}\PY{p}{)}    \PY{c+c1}{\PYZsh{} our new calculation for I0}
    \PY{n}{pred\PYZus{}vals\PYZus{}new} \PY{o}{=} \PY{n}{np}\PY{o}{.}\PY{n}{zeros}\PY{p}{(}\PY{n}{no\PYZus{}samp}\PY{p}{)} \PY{c+c1}{\PYZsh{} prediction for I, using the same formula as for I0, but with new values of k and M}

    \PY{c+c1}{\PYZsh{} Our parameter values.}
    \PY{n}{r\PYZus{}v} \PY{o}{=} \PY{n}{table}\PY{p}{[}\PY{l+s+s2}{\PYZdq{}}\PY{l+s+s2}{R}\PY{l+s+s2}{\PYZdq{}}\PY{p}{]}\PY{o}{.}\PY{n}{to\PYZus{}numpy}\PY{p}{(}\PY{p}{)}
    \PY{n}{b\PYZus{}v} \PY{o}{=} \PY{n}{table}\PY{p}{[}\PY{l+s+s2}{\PYZdq{}}\PY{l+s+s2}{B}\PY{l+s+s2}{\PYZdq{}}\PY{p}{]}\PY{o}{.}\PY{n}{to\PYZus{}numpy}\PY{p}{(}\PY{p}{)}
    \PY{n}{dk\PYZus{}v} \PY{o}{=} \PY{n}{table}\PY{p}{[}\PY{l+s+s2}{\PYZdq{}}\PY{l+s+s2}{dk}\PY{l+s+s2}{\PYZdq{}}\PY{p}{]}\PY{o}{.}\PY{n}{to\PYZus{}numpy}\PY{p}{(}\PY{p}{)}
    \PY{n}{dpsi0\PYZus{}v} \PY{o}{=} \PY{n}{table}\PY{p}{[}\PY{l+s+s2}{\PYZdq{}}\PY{l+s+s2}{dpsi0}\PY{l+s+s2}{\PYZdq{}}\PY{p}{]}\PY{o}{.}\PY{n}{to\PYZus{}numpy}\PY{p}{(}\PY{p}{)}
    \PY{n}{mu\PYZus{}v} \PY{o}{=} \PY{n}{table}\PY{p}{[}\PY{l+s+s2}{\PYZdq{}}\PY{l+s+s2}{mu}\PY{l+s+s2}{\PYZdq{}}\PY{p}{]}\PY{o}{.}\PY{n}{to\PYZus{}numpy}\PY{p}{(}\PY{p}{)}
    \PY{n}{dpsi\PYZus{}v} \PY{o}{=} \PY{n}{table}\PY{p}{[}\PY{l+s+s2}{\PYZdq{}}\PY{l+s+s2}{dpsi}\PY{l+s+s2}{\PYZdq{}}\PY{p}{]}\PY{o}{.}\PY{n}{to\PYZus{}numpy}\PY{p}{(}\PY{p}{)}

    \PY{c+c1}{\PYZsh{} Our M and k for the linear case.}
    \PY{n}{m\PYZus{}v} \PY{o}{=} \PY{n}{b\PYZus{}v} \PY{o}{*} \PY{n}{r\PYZus{}v} \PY{o}{*} \PY{n}{B\PYZus{}max} \PY{o}{*} \PY{n}{R\PYZus{}max} \PY{o}{*} \PY{n}{phi0inv}
    \PY{n}{k\PYZus{}v} \PY{o}{=} \PY{n}{k\PYZus{}val} \PY{o}{+} \PY{n}{dk\PYZus{}v} \PY{o}{/}\PY{l+m+mf}{2.0}\PY{o}{/}\PY{n}{R\PYZus{}max}

    \PY{c+c1}{\PYZsh{} Fill our storage arrays from the relevant functions.}
    \PY{k}{for} \PY{n}{i} \PY{o+ow}{in} \PY{n+nb}{range}\PY{p}{(}\PY{n}{no\PYZus{}samp}\PY{p}{)}\PY{p}{:}
        \PY{n}{pred\PYZus{}vals}\PY{p}{[}\PY{n}{i}\PY{p}{]} \PY{o}{=} \PY{n}{func\PYZus{}cur}\PY{p}{(}\PY{n}{r\PYZus{}v}\PY{p}{[}\PY{n}{i}\PY{p}{]}\PY{p}{,} \PY{n}{b\PYZus{}v}\PY{p}{[}\PY{n}{i}\PY{p}{]}\PY{p}{,} \PY{n}{k\PYZus{}val}\PY{p}{,} \PY{n}{dk\PYZus{}v}\PY{p}{[}\PY{n}{i}\PY{p}{]}\PY{p}{)}
    \PY{k}{for} \PY{n}{i} \PY{o+ow}{in} \PY{n+nb}{range}\PY{p}{(}\PY{n}{no\PYZus{}samp}\PY{p}{)}\PY{p}{:}
        \PY{n}{cur\PYZus{}vals}\PY{p}{[}\PY{n}{i}\PY{p}{]}\PY{o}{=}\PY{n}{func\PYZus{}cur\PYZus{}from\PYZus{}deriv}\PY{p}{(}\PY{n}{dpsi0\PYZus{}v}\PY{p}{[}\PY{n}{i}\PY{p}{]}\PY{p}{,} \PY{n}{m\PYZus{}v}\PY{p}{[}\PY{n}{i}\PY{p}{]}\PY{p}{,} \PY{n}{k\PYZus{}v}\PY{p}{[}\PY{n}{i}\PY{p}{]}\PY{p}{)}
    \PY{k}{for} \PY{n}{i} \PY{o+ow}{in} \PY{n+nb}{range}\PY{p}{(}\PY{n}{no\PYZus{}samp}\PY{p}{)}\PY{p}{:}
        \PY{n}{pred\PYZus{}vals\PYZus{}new}\PY{p}{[}\PY{n}{i}\PY{p}{]} \PY{o}{=} \PY{n}{func\PYZus{}cur\PYZus{}ee}\PY{p}{(}\PY{n}{r\PYZus{}v}\PY{p}{[}\PY{n}{i}\PY{p}{]}\PY{p}{,} \PY{n}{b\PYZus{}v}\PY{p}{[}\PY{n}{i}\PY{p}{]}\PY{p}{,} \PY{n}{dk\PYZus{}v}\PY{p}{[}\PY{n}{i}\PY{p}{]}\PY{p}{,} \PY{n}{mu\PYZus{}v}\PY{p}{[}\PY{n}{i}\PY{p}{]}\PY{p}{,} \PY{n}{dpsi\PYZus{}v}\PY{p}{[}\PY{n}{i}\PY{p}{]}\PY{p}{)}

    \PY{c+c1}{\PYZsh{} Add the columns to the table.}
    \PY{n}{table}\PY{p}{[}\PY{l+s+s2}{\PYZdq{}}\PY{l+s+s2}{I0 p1}\PY{l+s+s2}{\PYZdq{}}\PY{p}{]} \PY{o}{=} \PY{n}{cur\PYZus{}vals}
    \PY{n}{table}\PY{p}{[}\PY{l+s+s2}{\PYZdq{}}\PY{l+s+s2}{I0 p2}\PY{l+s+s2}{\PYZdq{}}\PY{p}{]} \PY{o}{=} \PY{n}{pred\PYZus{}vals}
    \PY{n}{table}\PY{p}{[}\PY{l+s+s2}{\PYZdq{}}\PY{l+s+s2}{I pred}\PY{l+s+s2}{\PYZdq{}}\PY{p}{]} \PY{o}{=} \PY{n}{pred\PYZus{}vals\PYZus{}new}
\end{Verbatim}
\end{tcolorbox}

    \begin{tcolorbox}[breakable, size=fbox, boxrule=1pt, pad at break*=1mm,colback=cellbackground, colframe=cellborder]
\prompt{In}{incolor}{37}{\boxspacing}
\begin{Verbatim}[commandchars=\\\{\}]
\PY{n}{predict\PYZus{}table}\PY{p}{(}\PY{n}{tbl021}\PY{p}{)}
\end{Verbatim}
\end{tcolorbox}

    \begin{tcolorbox}[breakable, size=fbox, boxrule=1pt, pad at break*=1mm,colback=cellbackground, colframe=cellborder]
\prompt{In}{incolor}{ }{\boxspacing}
\begin{Verbatim}[commandchars=\\\{\}]
\PY{c+c1}{\PYZsh{} Here are my triangle plots to compare the predicted and calculated values.}
\PY{c+c1}{\PYZsh{} I have a result of a line with a slope of 1, ie my guess is correct, not just proportional.}
\PY{c+c1}{\PYZsh{} There are a few outliers visible on the graph.}
\PY{n}{g} \PY{o}{=} \PY{n}{sns}\PY{o}{.}\PY{n}{PairGrid}\PY{p}{(}\PY{n}{tbl021}\PY{p}{[}\PY{p}{[}\PY{l+s+s2}{\PYZdq{}}\PY{l+s+s2}{R}\PY{l+s+s2}{\PYZdq{}}\PY{p}{,}\PY{l+s+s1}{\PYZsq{}}\PY{l+s+s1}{mu}\PY{l+s+s1}{\PYZsq{}}\PY{p}{,}\PY{l+s+s2}{\PYZdq{}}\PY{l+s+s2}{I v2}\PY{l+s+s2}{\PYZdq{}}\PY{p}{,}\PY{l+s+s2}{\PYZdq{}}\PY{l+s+s2}{I pred}\PY{l+s+s2}{\PYZdq{}}\PY{p}{]}\PY{p}{]}\PY{p}{)}

\PY{n}{g}\PY{o}{.}\PY{n}{map\PYZus{}diag}\PY{p}{(}\PY{n}{sns}\PY{o}{.}\PY{n}{kdeplot}\PY{p}{)}
\PY{n}{g}\PY{o}{.}\PY{n}{map\PYZus{}upper}\PY{p}{(}\PY{n}{sns}\PY{o}{.}\PY{n}{scatterplot}\PY{p}{)}
\PY{n}{g}\PY{o}{.}\PY{n}{map\PYZus{}lower}\PY{p}{(}\PY{n}{sns}\PY{o}{.}\PY{n}{scatterplot}\PY{p}{)}
\end{Verbatim}
\end{tcolorbox}

            \begin{tcolorbox}[breakable, size=fbox, boxrule=.5pt, pad at break*=1mm, opacityfill=0]
\prompt{Out}{outcolor}{ }{\boxspacing}
\begin{Verbatim}[commandchars=\\\{\}]
<seaborn.axisgrid.PairGrid at 0x27ba05c59a0>
\end{Verbatim}
\end{tcolorbox}
        
    \begin{center}
    \adjustimage{max size={0.9\linewidth}{0.9\paperheight}}{figs/bvp_fin_en/output_62_1.png}
    \end{center}
    { \hspace*{\fill} \\}
    
    And the model clearly shows not just a linear relationship, but rather
that we completely predicted the result. This should have been obvious
from the matching condition, but here it is clear. I will continue to
check these results, to show how good the results are. Here is our
measure of our fit.

    \begin{tcolorbox}[breakable, size=fbox, boxrule=1pt, pad at break*=1mm,colback=cellbackground, colframe=cellborder]
\prompt{In}{incolor}{39}{\boxspacing}
\begin{Verbatim}[commandchars=\\\{\}]
\PY{c+c1}{\PYZsh{} The variable names are from sklearn typical variable naming.}
\PY{c+c1}{\PYZsh{} since I am bad at naming variables.}

\PY{c+c1}{\PYZsh{} results from numerical calculation}
\PY{n}{y\PYZus{}test} \PY{o}{=} \PY{n}{tbl021}\PY{p}{[}\PY{l+s+s2}{\PYZdq{}}\PY{l+s+s2}{I v2}\PY{l+s+s2}{\PYZdq{}}\PY{p}{]}\PY{o}{.}\PY{n}{to\PYZus{}numpy}\PY{p}{(}\PY{p}{)}   
\PY{c+c1}{\PYZsh{} predicted model \PYZhy{}\PYZhy{} as the case without ee interaction, but with the new k and M values}
\PY{n}{y\PYZus{}pred} \PY{o}{=} \PY{n}{tbl021}\PY{p}{[}\PY{l+s+s2}{\PYZdq{}}\PY{l+s+s2}{I pred}\PY{l+s+s2}{\PYZdq{}}\PY{p}{]}\PY{o}{.}\PY{n}{to\PYZus{}numpy}\PY{p}{(}\PY{p}{)} 

\PY{c+c1}{\PYZsh{} For these estimates, I need enough decimal places for the error to be seen.}

\PY{n}{r2} \PY{o}{=} \PY{n}{r2\PYZus{}score}\PY{p}{(}\PY{n}{y\PYZus{}test}\PY{p}{,} \PY{n}{y\PYZus{}pred}\PY{p}{)}
\PY{n+nb}{print}\PY{p}{(}\PY{l+s+sa}{f}\PY{l+s+s2}{\PYZdq{}}\PY{l+s+s2}{R2:   }\PY{l+s+si}{\PYZob{}}\PY{n}{r2}\PY{l+s+si}{:}\PY{l+s+s2}{.7f}\PY{l+s+si}{\PYZcb{}}\PY{l+s+s2}{\PYZdq{}}\PY{p}{)}

\PY{n}{mse} \PY{o}{=} \PY{n}{mean\PYZus{}squared\PYZus{}error}\PY{p}{(}\PY{n}{y\PYZus{}test}\PY{p}{,} \PY{n}{y\PYZus{}pred}\PY{p}{)}
\PY{n+nb}{print}\PY{p}{(}\PY{l+s+sa}{f}\PY{l+s+s2}{\PYZdq{}}\PY{l+s+s2}{MSE:  }\PY{l+s+si}{\PYZob{}}\PY{n}{mse}\PY{l+s+si}{:}\PY{l+s+s2}{.7f}\PY{l+s+si}{\PYZcb{}}\PY{l+s+s2}{\PYZdq{}}\PY{p}{)}

\PY{n}{rmse} \PY{o}{=} \PY{n}{np}\PY{o}{.}\PY{n}{sqrt}\PY{p}{(}\PY{n}{mse}\PY{p}{)}
\PY{n+nb}{print}\PY{p}{(}\PY{l+s+sa}{f}\PY{l+s+s2}{\PYZdq{}}\PY{l+s+s2}{RMSE: }\PY{l+s+si}{\PYZob{}}\PY{n}{rmse}\PY{l+s+si}{:}\PY{l+s+s2}{.7f}\PY{l+s+si}{\PYZcb{}}\PY{l+s+s2}{\PYZdq{}}\PY{p}{)}

\PY{n}{mae} \PY{o}{=} \PY{n}{mean\PYZus{}absolute\PYZus{}error}\PY{p}{(}\PY{n}{y\PYZus{}test}\PY{p}{,} \PY{n}{y\PYZus{}pred}\PY{p}{)}
\PY{n+nb}{print}\PY{p}{(}\PY{l+s+sa}{f}\PY{l+s+s2}{\PYZdq{}}\PY{l+s+s2}{MAE:  }\PY{l+s+si}{\PYZob{}}\PY{n}{mae}\PY{l+s+si}{:}\PY{l+s+s2}{.7f}\PY{l+s+si}{\PYZcb{}}\PY{l+s+s2}{\PYZdq{}}\PY{p}{)}

\PY{c+c1}{\PYZsh{} This is not the correct metric for this measure, but still gives a decent}
\PY{c+c1}{\PYZsh{} understanding of how far the points can diverge from the plot. }
\PY{c+c1}{\PYZsh{} Based on the other measures, these points are most likely erroneous.}
\PY{c+c1}{\PYZsh{} I know the error in the calculated current can be large.}
\PY{n}{max\PYZus{}val} \PY{o}{=} \PY{n}{np}\PY{o}{.}\PY{n}{max}\PY{p}{(}\PY{p}{[}\PY{n}{np}\PY{o}{.}\PY{n}{max}\PY{p}{(}\PY{n}{y\PYZus{}test}\PY{p}{)}\PY{p}{,} \PY{n}{np}\PY{o}{.}\PY{n}{max}\PY{p}{(}\PY{n}{y\PYZus{}pred}\PY{p}{)}\PY{p}{]}\PY{p}{)}
\PY{n}{max\PYZus{}sep} \PY{o}{=} \PY{n}{np}\PY{o}{.}\PY{n}{max}\PY{p}{(}\PY{n}{y\PYZus{}test}\PY{o}{\PYZhy{}}\PY{n}{y\PYZus{}pred}\PY{p}{)}
\PY{n+nb}{print}\PY{p}{(}\PY{l+s+sa}{f}\PY{l+s+s2}{\PYZdq{}}\PY{l+s+s2}{Estimate maximum error: }\PY{l+s+si}{\PYZob{}}\PY{n}{max\PYZus{}sep}\PY{o}{/}\PY{n}{max\PYZus{}val}\PY{l+s+si}{\PYZcb{}}\PY{l+s+s2}{\PYZdq{}}\PY{p}{)}
\end{Verbatim}
\end{tcolorbox}

    \begin{Verbatim}[commandchars=\\\{\}]
R2:   0.9998840
MSE:  0.0000001
RMSE: 0.0003233
MAE:  0.0000038
Estimate maximum error: 0.2835049216758385
    \end{Verbatim}

    Now lets try for our other datasets.

    \begin{tcolorbox}[breakable, size=fbox, boxrule=1pt, pad at break*=1mm,colback=cellbackground, colframe=cellborder]
\prompt{In}{incolor}{40}{\boxspacing}
\begin{Verbatim}[commandchars=\\\{\}]
\PY{n}{predict\PYZus{}table}\PY{p}{(}\PY{n}{tbl023}\PY{p}{)}
\end{Verbatim}
\end{tcolorbox}

    \begin{tcolorbox}[breakable, size=fbox, boxrule=1pt, pad at break*=1mm,colback=cellbackground, colframe=cellborder]
\prompt{In}{incolor}{41}{\boxspacing}
\begin{Verbatim}[commandchars=\\\{\}]
\PY{n}{tbl023}\PY{o}{.}\PY{n}{head}\PY{p}{(}\PY{p}{)}
\end{Verbatim}
\end{tcolorbox}

            \begin{tcolorbox}[breakable, size=fbox, boxrule=.5pt, pad at break*=1mm, opacityfill=0]
\prompt{Out}{outcolor}{41}{\boxspacing}
\begin{Verbatim}[commandchars=\\\{\}]
      R     B   dk            mu  \textbackslash{}
0  0.57  0.11  0.1  1.000000e-10
1  0.57  0.11  0.1  1.584893e-10
2  0.57  0.11  0.1  2.511886e-10
3  0.57  0.11  0.1  3.981072e-10
4  0.57  0.11  0.1  6.309573e-10

                                            dpsi0                  a0  \textbackslash{}
0 -1.068926e+11-5.376625e+                    10j -1.740277+4.453840j
1 -1.068926e+11-5.376625e+                    10j -1.740277+4.453840j
2 -1.068926e+11-5.376625e+                    10j -1.740277+4.453840j
3 -1.068926e+11-5.376625e+                    10j -1.740277+4.453840j
4 -1.068926e+11-5.376625e+                    10j -1.740277+4.453840j

                   b0    A max 0        I0  \textbackslash{}
0  2.740277-4.453840j  10.011085 -4.480554
1  2.740277-4.453840j  10.011085 -4.480554
2  2.740277-4.453840j  10.011085 -4.480554
3  2.740277-4.453840j  10.011085 -4.480554
4  2.740277-4.453840j  10.011085 -4.480554

                                             dpsi                   a  \textbackslash{}
0  2.249675e+11+1.115203e+                    11j  5.146633-9.373606j
1  1.847757e+11+9.159280e+                    10j  4.316325-7.698956j
2  1.522990e+11+7.549034e+                    10j  3.645392-6.345764j
3  1.259541e+11+6.243197e+                    10j  3.101296-5.248064j
4  1.044921e+11+5.180113e+                    10j  2.658346-4.353821j

                    b  A max new      I v2  effective mu     I0 p1     I0 p2  \textbackslash{}
0 -4.146633+9.373606j  20.943396  9.293266  4.386258e-08 -4.480554 -4.480554
1 -3.316325+7.698956j  17.209193  7.632650  4.693761e-08 -4.480554 -4.480554
2 -2.645392+6.345764j  14.193394  6.290784  5.060256e-08 -4.480554 -4.480554
3 -2.101296+5.248064j  11.749025  5.202591  5.495455e-08 -4.480554 -4.480554
4 -1.658346+4.353821j   9.760186  4.316691  6.010578e-08 -4.480554 -4.480554

     I pred
0  9.293387
1  7.632756
2  6.290878
3  5.202676
4  4.316768
\end{Verbatim}
\end{tcolorbox}
        
    \begin{tcolorbox}[breakable, size=fbox, boxrule=1pt, pad at break*=1mm,colback=cellbackground, colframe=cellborder]
\prompt{In}{incolor}{42}{\boxspacing}
\begin{Verbatim}[commandchars=\\\{\}]
\PY{n}{g} \PY{o}{=} \PY{n}{sns}\PY{o}{.}\PY{n}{PairGrid}\PY{p}{(}\PY{n}{tbl023}\PY{p}{[}\PY{p}{[}\PY{l+s+s2}{\PYZdq{}}\PY{l+s+s2}{R}\PY{l+s+s2}{\PYZdq{}}\PY{p}{,}\PY{l+s+s1}{\PYZsq{}}\PY{l+s+s1}{mu}\PY{l+s+s1}{\PYZsq{}}\PY{p}{,}\PY{l+s+s2}{\PYZdq{}}\PY{l+s+s2}{I v2}\PY{l+s+s2}{\PYZdq{}}\PY{p}{,}\PY{l+s+s2}{\PYZdq{}}\PY{l+s+s2}{I pred}\PY{l+s+s2}{\PYZdq{}}\PY{p}{]}\PY{p}{]}\PY{p}{)}

\PY{n}{g}\PY{o}{.}\PY{n}{map\PYZus{}diag}\PY{p}{(}\PY{n}{sns}\PY{o}{.}\PY{n}{kdeplot}\PY{p}{)}
\PY{n}{g}\PY{o}{.}\PY{n}{map\PYZus{}upper}\PY{p}{(}\PY{n}{sns}\PY{o}{.}\PY{n}{scatterplot}\PY{p}{)}
\PY{n}{g}\PY{o}{.}\PY{n}{map\PYZus{}lower}\PY{p}{(}\PY{n}{sns}\PY{o}{.}\PY{n}{scatterplot}\PY{p}{)}
\end{Verbatim}
\end{tcolorbox}

            \begin{tcolorbox}[breakable, size=fbox, boxrule=.5pt, pad at break*=1mm, opacityfill=0]
\prompt{Out}{outcolor}{42}{\boxspacing}
\begin{Verbatim}[commandchars=\\\{\}]
<seaborn.axisgrid.PairGrid at 0x27b9f53d820>
\end{Verbatim}
\end{tcolorbox}
        
    \begin{center}
    \adjustimage{max size={0.9\linewidth}{0.9\paperheight}}{figs/bvp_fin_en/output_68_1.png}
    \end{center}
    { \hspace*{\fill} \\}
    
    And for testing our goodness of fit.

    \begin{tcolorbox}[breakable, size=fbox, boxrule=1pt, pad at break*=1mm,colback=cellbackground, colframe=cellborder]
\prompt{In}{incolor}{43}{\boxspacing}
\begin{Verbatim}[commandchars=\\\{\}]
\PY{c+c1}{\PYZsh{} The variable names are from sklearn typical variable naming.}
\PY{c+c1}{\PYZsh{} since I am bad at naming variables.}

\PY{c+c1}{\PYZsh{} results from numerical calculation}
\PY{n}{y\PYZus{}test} \PY{o}{=} \PY{n}{tbl023}\PY{p}{[}\PY{l+s+s2}{\PYZdq{}}\PY{l+s+s2}{I v2}\PY{l+s+s2}{\PYZdq{}}\PY{p}{]}\PY{o}{.}\PY{n}{to\PYZus{}numpy}\PY{p}{(}\PY{p}{)}   
\PY{c+c1}{\PYZsh{} predicted model \PYZhy{}\PYZhy{} as the case without ee interaction, but with the new k and M values}
\PY{n}{y\PYZus{}pred} \PY{o}{=} \PY{n}{tbl023}\PY{p}{[}\PY{l+s+s2}{\PYZdq{}}\PY{l+s+s2}{I pred}\PY{l+s+s2}{\PYZdq{}}\PY{p}{]}\PY{o}{.}\PY{n}{to\PYZus{}numpy}\PY{p}{(}\PY{p}{)} 

\PY{c+c1}{\PYZsh{} For these estimates, I need enough decimal places for the error to be seen.}

\PY{n}{r2} \PY{o}{=} \PY{n}{r2\PYZus{}score}\PY{p}{(}\PY{n}{y\PYZus{}test}\PY{p}{,} \PY{n}{y\PYZus{}pred}\PY{p}{)}
\PY{n+nb}{print}\PY{p}{(}\PY{l+s+sa}{f}\PY{l+s+s2}{\PYZdq{}}\PY{l+s+s2}{R2:   }\PY{l+s+si}{\PYZob{}}\PY{n}{r2}\PY{l+s+si}{:}\PY{l+s+s2}{.7f}\PY{l+s+si}{\PYZcb{}}\PY{l+s+s2}{\PYZdq{}}\PY{p}{)}

\PY{n}{mse} \PY{o}{=} \PY{n}{mean\PYZus{}squared\PYZus{}error}\PY{p}{(}\PY{n}{y\PYZus{}test}\PY{p}{,} \PY{n}{y\PYZus{}pred}\PY{p}{)}
\PY{n+nb}{print}\PY{p}{(}\PY{l+s+sa}{f}\PY{l+s+s2}{\PYZdq{}}\PY{l+s+s2}{MSE:  }\PY{l+s+si}{\PYZob{}}\PY{n}{mse}\PY{l+s+si}{:}\PY{l+s+s2}{.7f}\PY{l+s+si}{\PYZcb{}}\PY{l+s+s2}{\PYZdq{}}\PY{p}{)}

\PY{n}{rmse} \PY{o}{=} \PY{n}{np}\PY{o}{.}\PY{n}{sqrt}\PY{p}{(}\PY{n}{mse}\PY{p}{)}
\PY{n+nb}{print}\PY{p}{(}\PY{l+s+sa}{f}\PY{l+s+s2}{\PYZdq{}}\PY{l+s+s2}{RMSE: }\PY{l+s+si}{\PYZob{}}\PY{n}{rmse}\PY{l+s+si}{:}\PY{l+s+s2}{.7f}\PY{l+s+si}{\PYZcb{}}\PY{l+s+s2}{\PYZdq{}}\PY{p}{)}

\PY{n}{mae} \PY{o}{=} \PY{n}{mean\PYZus{}absolute\PYZus{}error}\PY{p}{(}\PY{n}{y\PYZus{}test}\PY{p}{,} \PY{n}{y\PYZus{}pred}\PY{p}{)}
\PY{n+nb}{print}\PY{p}{(}\PY{l+s+sa}{f}\PY{l+s+s2}{\PYZdq{}}\PY{l+s+s2}{MAE:  }\PY{l+s+si}{\PYZob{}}\PY{n}{mae}\PY{l+s+si}{:}\PY{l+s+s2}{.7f}\PY{l+s+si}{\PYZcb{}}\PY{l+s+s2}{\PYZdq{}}\PY{p}{)}

\PY{c+c1}{\PYZsh{} This is not the correct metric for this measure, but still gives a decent}
\PY{c+c1}{\PYZsh{} understanding of how far the points can diverge from the plot. }
\PY{c+c1}{\PYZsh{} Based on the other measures, these points are most likely erroneous.}
\PY{c+c1}{\PYZsh{} I know the error in the calculated current can be large.}
\PY{n}{max\PYZus{}val} \PY{o}{=} \PY{n}{np}\PY{o}{.}\PY{n}{max}\PY{p}{(}\PY{p}{[}\PY{n}{np}\PY{o}{.}\PY{n}{max}\PY{p}{(}\PY{n}{y\PYZus{}test}\PY{p}{)}\PY{p}{,} \PY{n}{np}\PY{o}{.}\PY{n}{max}\PY{p}{(}\PY{n}{y\PYZus{}pred}\PY{p}{)}\PY{p}{]}\PY{p}{)}
\PY{n}{max\PYZus{}sep} \PY{o}{=} \PY{n}{np}\PY{o}{.}\PY{n}{max}\PY{p}{(}\PY{n}{y\PYZus{}test}\PY{o}{\PYZhy{}}\PY{n}{y\PYZus{}pred}\PY{p}{)}
\PY{n+nb}{print}\PY{p}{(}\PY{l+s+sa}{f}\PY{l+s+s2}{\PYZdq{}}\PY{l+s+s2}{Estimate maximum error: }\PY{l+s+si}{\PYZob{}}\PY{n}{max\PYZus{}sep}\PY{o}{/}\PY{n}{max\PYZus{}val}\PY{l+s+si}{\PYZcb{}}\PY{l+s+s2}{\PYZdq{}}\PY{p}{)}
\end{Verbatim}
\end{tcolorbox}

    \begin{Verbatim}[commandchars=\\\{\}]
R2:   0.9999993
MSE:  0.0000019
RMSE: 0.0013798
MAE:  0.0001261
Estimate maximum error: 0.0030560971707368875
    \end{Verbatim}

    And for the Monte Carlo dataset.

    \begin{tcolorbox}[breakable, size=fbox, boxrule=1pt, pad at break*=1mm,colback=cellbackground, colframe=cellborder]
\prompt{In}{incolor}{44}{\boxspacing}
\begin{Verbatim}[commandchars=\\\{\}]
\PY{n}{df} \PY{o}{=} \PY{n}{tbl025}
\PY{n}{predict\PYZus{}table}\PY{p}{(}\PY{n}{df}\PY{p}{)}
\end{Verbatim}
\end{tcolorbox}

    \begin{tcolorbox}[breakable, size=fbox, boxrule=1pt, pad at break*=1mm,colback=cellbackground, colframe=cellborder]
\prompt{In}{incolor}{45}{\boxspacing}
\begin{Verbatim}[commandchars=\\\{\}]
\PY{n}{g} \PY{o}{=} \PY{n}{sns}\PY{o}{.}\PY{n}{PairGrid}\PY{p}{(}\PY{n}{df}\PY{p}{[}\PY{p}{[}\PY{l+s+s2}{\PYZdq{}}\PY{l+s+s2}{R}\PY{l+s+s2}{\PYZdq{}}\PY{p}{,}\PY{l+s+s1}{\PYZsq{}}\PY{l+s+s1}{mu}\PY{l+s+s1}{\PYZsq{}}\PY{p}{,}\PY{l+s+s2}{\PYZdq{}}\PY{l+s+s2}{I v2}\PY{l+s+s2}{\PYZdq{}}\PY{p}{,}\PY{l+s+s2}{\PYZdq{}}\PY{l+s+s2}{I pred}\PY{l+s+s2}{\PYZdq{}}\PY{p}{]}\PY{p}{]}\PY{p}{)}

\PY{n}{g}\PY{o}{.}\PY{n}{map\PYZus{}diag}\PY{p}{(}\PY{n}{sns}\PY{o}{.}\PY{n}{kdeplot}\PY{p}{)}
\PY{n}{g}\PY{o}{.}\PY{n}{map\PYZus{}upper}\PY{p}{(}\PY{n}{sns}\PY{o}{.}\PY{n}{scatterplot}\PY{p}{)}
\PY{n}{g}\PY{o}{.}\PY{n}{map\PYZus{}lower}\PY{p}{(}\PY{n}{sns}\PY{o}{.}\PY{n}{scatterplot}\PY{p}{)}
\end{Verbatim}
\end{tcolorbox}

            \begin{tcolorbox}[breakable, size=fbox, boxrule=.5pt, pad at break*=1mm, opacityfill=0]
\prompt{Out}{outcolor}{45}{\boxspacing}
\begin{Verbatim}[commandchars=\\\{\}]
<seaborn.axisgrid.PairGrid at 0x27ba34166f0>
\end{Verbatim}
\end{tcolorbox}
        
    \begin{center}
    \adjustimage{max size={0.9\linewidth}{0.9\paperheight}}{figs/bvp_fin_en/output_73_1.png}
    \end{center}
    { \hspace*{\fill} \\}
    
    \begin{tcolorbox}[breakable, size=fbox, boxrule=1pt, pad at break*=1mm,colback=cellbackground, colframe=cellborder]
\prompt{In}{incolor}{46}{\boxspacing}
\begin{Verbatim}[commandchars=\\\{\}]
\PY{c+c1}{\PYZsh{} The variable names are from sklearn typical variable naming.}
\PY{c+c1}{\PYZsh{} since I am bad at naming variables.}

\PY{c+c1}{\PYZsh{} results from numerical calculation}
\PY{n}{y\PYZus{}test} \PY{o}{=} \PY{n}{df}\PY{p}{[}\PY{l+s+s2}{\PYZdq{}}\PY{l+s+s2}{I v2}\PY{l+s+s2}{\PYZdq{}}\PY{p}{]}\PY{o}{.}\PY{n}{to\PYZus{}numpy}\PY{p}{(}\PY{p}{)}   
\PY{c+c1}{\PYZsh{} predicted model \PYZhy{}\PYZhy{} as the case without ee interaction, but with the new k and M values}
\PY{n}{y\PYZus{}pred} \PY{o}{=} \PY{n}{df}\PY{p}{[}\PY{l+s+s2}{\PYZdq{}}\PY{l+s+s2}{I pred}\PY{l+s+s2}{\PYZdq{}}\PY{p}{]}\PY{o}{.}\PY{n}{to\PYZus{}numpy}\PY{p}{(}\PY{p}{)} 

\PY{c+c1}{\PYZsh{} For these estimates, I need enough decimal places for the error to be seen.}

\PY{n}{r2} \PY{o}{=} \PY{n}{r2\PYZus{}score}\PY{p}{(}\PY{n}{y\PYZus{}test}\PY{p}{,} \PY{n}{y\PYZus{}pred}\PY{p}{)}
\PY{n+nb}{print}\PY{p}{(}\PY{l+s+sa}{f}\PY{l+s+s2}{\PYZdq{}}\PY{l+s+s2}{R2:   }\PY{l+s+si}{\PYZob{}}\PY{n}{r2}\PY{l+s+si}{:}\PY{l+s+s2}{.7f}\PY{l+s+si}{\PYZcb{}}\PY{l+s+s2}{\PYZdq{}}\PY{p}{)}

\PY{n}{mse} \PY{o}{=} \PY{n}{mean\PYZus{}squared\PYZus{}error}\PY{p}{(}\PY{n}{y\PYZus{}test}\PY{p}{,} \PY{n}{y\PYZus{}pred}\PY{p}{)}
\PY{n+nb}{print}\PY{p}{(}\PY{l+s+sa}{f}\PY{l+s+s2}{\PYZdq{}}\PY{l+s+s2}{MSE:  }\PY{l+s+si}{\PYZob{}}\PY{n}{mse}\PY{l+s+si}{:}\PY{l+s+s2}{.7f}\PY{l+s+si}{\PYZcb{}}\PY{l+s+s2}{\PYZdq{}}\PY{p}{)}

\PY{n}{rmse} \PY{o}{=} \PY{n}{np}\PY{o}{.}\PY{n}{sqrt}\PY{p}{(}\PY{n}{mse}\PY{p}{)}
\PY{n+nb}{print}\PY{p}{(}\PY{l+s+sa}{f}\PY{l+s+s2}{\PYZdq{}}\PY{l+s+s2}{RMSE: }\PY{l+s+si}{\PYZob{}}\PY{n}{rmse}\PY{l+s+si}{:}\PY{l+s+s2}{.7f}\PY{l+s+si}{\PYZcb{}}\PY{l+s+s2}{\PYZdq{}}\PY{p}{)}

\PY{n}{mae} \PY{o}{=} \PY{n}{mean\PYZus{}absolute\PYZus{}error}\PY{p}{(}\PY{n}{y\PYZus{}test}\PY{p}{,} \PY{n}{y\PYZus{}pred}\PY{p}{)}
\PY{n+nb}{print}\PY{p}{(}\PY{l+s+sa}{f}\PY{l+s+s2}{\PYZdq{}}\PY{l+s+s2}{MAE:  }\PY{l+s+si}{\PYZob{}}\PY{n}{mae}\PY{l+s+si}{:}\PY{l+s+s2}{.7f}\PY{l+s+si}{\PYZcb{}}\PY{l+s+s2}{\PYZdq{}}\PY{p}{)}

\PY{c+c1}{\PYZsh{} This is not the correct metric for this measure, but still gives a decent}
\PY{c+c1}{\PYZsh{} understanding of how far the points can diverge from the plot. }
\PY{c+c1}{\PYZsh{} Based on the other measures, these points are most likely erroneous.}
\PY{c+c1}{\PYZsh{} I know the error in the calculated current can be large.}
\PY{n}{max\PYZus{}val} \PY{o}{=} \PY{n}{np}\PY{o}{.}\PY{n}{max}\PY{p}{(}\PY{p}{[}\PY{n}{np}\PY{o}{.}\PY{n}{max}\PY{p}{(}\PY{n}{y\PYZus{}test}\PY{p}{)}\PY{p}{,} \PY{n}{np}\PY{o}{.}\PY{n}{max}\PY{p}{(}\PY{n}{y\PYZus{}pred}\PY{p}{)}\PY{p}{]}\PY{p}{)}
\PY{n}{max\PYZus{}sep} \PY{o}{=} \PY{n}{np}\PY{o}{.}\PY{n}{max}\PY{p}{(}\PY{n}{y\PYZus{}test}\PY{o}{\PYZhy{}}\PY{n}{y\PYZus{}pred}\PY{p}{)}
\PY{n+nb}{print}\PY{p}{(}\PY{l+s+sa}{f}\PY{l+s+s2}{\PYZdq{}}\PY{l+s+s2}{Estimate maximum error: }\PY{l+s+si}{\PYZob{}}\PY{n}{max\PYZus{}sep}\PY{o}{/}\PY{n}{max\PYZus{}val}\PY{l+s+si}{\PYZcb{}}\PY{l+s+s2}{\PYZdq{}}\PY{p}{)}
\end{Verbatim}
\end{tcolorbox}

    \begin{Verbatim}[commandchars=\\\{\}]
R2:   1.0000000
MSE:  0.0000000
RMSE: 0.0000002
MAE:  0.0000001
Estimate maximum error: 0.00038383139285069875
    \end{Verbatim}

    This is definitely the correct equation.

    \hypertarget{test-a-b}{%
\section{Test a, b}\label{test-a-b}}

    The last analysis of this notebook is to check that \(a_j\) and \(b_j\)
have the same form as the linear case, just with new \(k_j\) and \(M_j\)
values.

    \begin{tcolorbox}[breakable, size=fbox, boxrule=1pt, pad at break*=1mm,colback=cellbackground, colframe=cellborder]
\prompt{In}{incolor}{47}{\boxspacing}
\begin{Verbatim}[commandchars=\\\{\}]
\PY{c+c1}{\PYZsh{} The linear model for a and b.}
\PY{c+c1}{\PYZsh{} Requires k, M = BR pi / Phi\PYZus{}0, and R.}
\PY{c+c1}{\PYZsh{} Returns a\PYZus{}j, b\PYZus{}j.}
\PY{k}{def} \PY{n+nf}{ab\PYZus{}pred}\PY{p}{(}\PY{n}{k}\PY{p}{,} \PY{n}{M}\PY{p}{,} \PY{n}{R}\PY{p}{)}\PY{p}{:}
    \PY{n}{denom} \PY{o}{=} \PY{l+m+mi}{2}\PY{n}{j}\PY{o}{*}\PY{n}{np}\PY{o}{.}\PY{n}{sin}\PY{p}{(}\PY{l+m+mi}{2}\PY{o}{*}\PY{n}{pi}\PY{o}{*}\PY{n}{k}\PY{o}{*}\PY{n}{R}\PY{p}{)}\PY{o}{*}\PY{n}{np}\PY{o}{.}\PY{n}{exp}\PY{p}{(}\PY{l+m+mi}{2}\PY{n}{j}\PY{o}{*}\PY{n}{M}\PY{o}{*}\PY{l+m+mi}{2}\PY{o}{*}\PY{n}{pi}\PY{o}{*}\PY{n}{R}\PY{p}{)}
    \PY{n}{a\PYZus{}num} \PY{o}{=} \PY{l+m+mi}{1} \PY{o}{\PYZhy{}} \PY{n}{np}\PY{o}{.}\PY{n}{exp}\PY{p}{(}\PY{l+m+mi}{1}\PY{n}{j}\PY{o}{*}\PY{l+m+mi}{2}\PY{o}{*}\PY{n}{pi}\PY{o}{*}\PY{n}{R}\PY{o}{*}\PY{p}{(}\PY{o}{\PYZhy{}}\PY{n}{k}\PY{o}{+}\PY{n}{M}\PY{p}{)}\PY{p}{)}
    \PY{n}{b\PYZus{}num} \PY{o}{=} \PY{o}{\PYZhy{}}\PY{l+m+mi}{1} \PY{o}{+} \PY{n}{np}\PY{o}{.}\PY{n}{exp}\PY{p}{(}\PY{l+m+mi}{1}\PY{n}{j}\PY{o}{*}\PY{l+m+mi}{2}\PY{o}{*}\PY{n}{pi}\PY{o}{*}\PY{n}{R}\PY{o}{*}\PY{p}{(}\PY{n}{k}\PY{o}{+}\PY{n}{M}\PY{p}{)}\PY{p}{)}
    \PY{k}{return} \PY{n}{a\PYZus{}num} \PY{o}{/} \PY{n}{denom}\PY{p}{,} \PY{n}{b\PYZus{}num} \PY{o}{/} \PY{n}{denom}
\end{Verbatim}
\end{tcolorbox}

    \begin{tcolorbox}[breakable, size=fbox, boxrule=1pt, pad at break*=1mm,colback=cellbackground, colframe=cellborder]
\prompt{In}{incolor}{48}{\boxspacing}
\begin{Verbatim}[commandchars=\\\{\}]
\PY{c+c1}{\PYZsh{} The nonlinear model for a and b.}
\PY{c+c1}{\PYZsh{} It is the same as the linear model, but with new k\PYZus{}j and M\PYZus{}j values.}
\PY{c+c1}{\PYZsh{} Requires R, B, dk, mu, dpsi0, A (1 by default), and k (k\PYZus{}Fermi for gold by default)}
\PY{c+c1}{\PYZsh{} Returns a\PYZus{}j, b\PYZus{}j.}
\PY{k}{def} \PY{n+nf}{ab\PYZus{}pred\PYZus{}ee}\PY{p}{(}\PY{n}{R}\PY{p}{,} \PY{n}{B}\PY{p}{,} \PY{n}{dk}\PY{p}{,} \PY{n}{mu}\PY{p}{,} \PY{n}{dpsi0}\PY{p}{,} \PY{n}{A} \PY{o}{=} \PY{l+m+mf}{1.}\PY{p}{,} \PY{n}{k} \PY{o}{=} \PY{n}{eelib}\PY{o}{.}\PY{n}{kFAu}\PY{p}{)}\PY{p}{:}
    \PY{n}{R} \PY{o}{=} \PY{n}{R} \PY{o}{*} \PY{n}{R\PYZus{}max} \PY{c+c1}{\PYZsh{} rescale R}
    \PY{n}{B} \PY{o}{=} \PY{n}{B} \PY{o}{*} \PY{n}{B\PYZus{}max} \PY{c+c1}{\PYZsh{} rescale B}
    \PY{n}{k\PYZus{}new} \PY{o}{=} \PY{n}{eelib}\PY{o}{.}\PY{n}{k\PYZus{}M\PYZus{}models\PYZus{}ivp}\PY{o}{.}\PY{n}{pred\PYZus{}fast\PYZus{}k\PYZus{}true}\PY{p}{(}\PY{n}{dpsi0}\PY{p}{,} \PY{n}{mu}\PY{p}{,} \PY{n}{dk}\PY{p}{,} \PY{n}{B}\PY{p}{,} \PY{n}{R}\PY{p}{,} \PY{n}{A}\PY{o}{=}\PY{n}{A}\PY{p}{,} \PY{n}{k0}\PY{o}{=}\PY{n}{k}\PY{p}{)} \PY{c+c1}{\PYZsh{} predict k}
    \PY{n}{M\PYZus{}new} \PY{o}{=} \PY{n}{eelib}\PY{o}{.}\PY{n}{pred\PYZus{}slow\PYZus{}k\PYZus{}v3}\PY{p}{(}\PY{n}{dpsi0}\PY{p}{,} \PY{n}{mu}\PY{p}{,} \PY{n}{dk}\PY{p}{,} \PY{n}{B}\PY{p}{,} \PY{n}{R}\PY{p}{,} \PY{n}{A}\PY{o}{=}\PY{n}{A}\PY{p}{,} \PY{n}{k0}\PY{o}{=}\PY{n}{k}\PY{p}{)} \PY{c+c1}{\PYZsh{} predict M}
    
    \PY{n}{ret\PYZus{}val1}\PY{p}{,} \PY{n}{ret\PYZus{}val2} \PY{o}{=} \PY{n}{ab\PYZus{}pred}\PY{p}{(}\PY{n}{k\PYZus{}new}\PY{p}{,} \PY{n}{M\PYZus{}new}\PY{p}{,} \PY{n}{R}\PY{p}{)} \PY{c+c1}{\PYZsh{} as the linear case, but with new values of k and M}

    \PY{k}{return} \PY{n}{ret\PYZus{}val1}\PY{p}{,} \PY{n}{ret\PYZus{}val2}
\end{Verbatim}
\end{tcolorbox}

    \begin{tcolorbox}[breakable, size=fbox, boxrule=1pt, pad at break*=1mm,colback=cellbackground, colframe=cellborder]
\prompt{In}{incolor}{49}{\boxspacing}
\begin{Verbatim}[commandchars=\\\{\}]
\PY{c+c1}{\PYZsh{} Add the predicted nonlinear a\PYZus{}j and b\PYZus{}j values to the table}
\PY{k}{def} \PY{n+nf}{predict\PYZus{}table\PYZus{}ab}\PY{p}{(}\PY{n}{table}\PY{p}{,} \PY{n}{k\PYZus{}val}\PY{o}{=}\PY{n}{eelib}\PY{o}{.}\PY{n}{kFAu}\PY{p}{,} \PY{n}{A\PYZus{}val}\PY{o}{=}\PY{l+m+mf}{1.0}\PY{p}{)}\PY{p}{:}

    \PY{n}{no\PYZus{}samp} \PY{o}{=} \PY{n+nb}{len}\PY{p}{(}\PY{n}{table}\PY{o}{.}\PY{n}{index}\PY{p}{)} \PY{c+c1}{\PYZsh{} how many rows}

    \PY{n}{pred\PYZus{}vals\PYZus{}a} \PY{o}{=} \PY{n}{np}\PY{o}{.}\PY{n}{zeros}\PY{p}{(}\PY{n}{no\PYZus{}samp}\PY{p}{,} \PY{n+nb}{complex}\PY{p}{)} \PY{c+c1}{\PYZsh{} preallocate space}
    \PY{n}{pred\PYZus{}vals\PYZus{}b} \PY{o}{=} \PY{n}{np}\PY{o}{.}\PY{n}{zeros}\PY{p}{(}\PY{n}{no\PYZus{}samp}\PY{p}{,} \PY{n+nb}{complex}\PY{p}{)} \PY{c+c1}{\PYZsh{} preallocate space}

    \PY{c+c1}{\PYZsh{} our columns of parameters}
    \PY{n}{r\PYZus{}v} \PY{o}{=} \PY{n}{table}\PY{p}{[}\PY{l+s+s2}{\PYZdq{}}\PY{l+s+s2}{R}\PY{l+s+s2}{\PYZdq{}}\PY{p}{]}\PY{o}{.}\PY{n}{to\PYZus{}numpy}\PY{p}{(}\PY{p}{)}
    \PY{n}{b\PYZus{}v} \PY{o}{=} \PY{n}{table}\PY{p}{[}\PY{l+s+s2}{\PYZdq{}}\PY{l+s+s2}{B}\PY{l+s+s2}{\PYZdq{}}\PY{p}{]}\PY{o}{.}\PY{n}{to\PYZus{}numpy}\PY{p}{(}\PY{p}{)}
    \PY{n}{dk\PYZus{}v} \PY{o}{=} \PY{n}{table}\PY{p}{[}\PY{l+s+s2}{\PYZdq{}}\PY{l+s+s2}{dk}\PY{l+s+s2}{\PYZdq{}}\PY{p}{]}\PY{o}{.}\PY{n}{to\PYZus{}numpy}\PY{p}{(}\PY{p}{)}
    \PY{n}{dpsi0\PYZus{}v} \PY{o}{=} \PY{n}{table}\PY{p}{[}\PY{l+s+s2}{\PYZdq{}}\PY{l+s+s2}{dpsi0}\PY{l+s+s2}{\PYZdq{}}\PY{p}{]}\PY{o}{.}\PY{n}{to\PYZus{}numpy}\PY{p}{(}\PY{p}{)}
    \PY{n}{mu\PYZus{}v} \PY{o}{=} \PY{n}{table}\PY{p}{[}\PY{l+s+s2}{\PYZdq{}}\PY{l+s+s2}{mu}\PY{l+s+s2}{\PYZdq{}}\PY{p}{]}\PY{o}{.}\PY{n}{to\PYZus{}numpy}\PY{p}{(}\PY{p}{)}
    \PY{n}{dpsi\PYZus{}v} \PY{o}{=} \PY{n}{table}\PY{p}{[}\PY{l+s+s2}{\PYZdq{}}\PY{l+s+s2}{dpsi}\PY{l+s+s2}{\PYZdq{}}\PY{p}{]}\PY{o}{.}\PY{n}{to\PYZus{}numpy}\PY{p}{(}\PY{p}{)}

    \PY{c+c1}{\PYZsh{} linear model k and M}
    \PY{n}{m\PYZus{}v} \PY{o}{=} \PY{n}{b\PYZus{}v} \PY{o}{*} \PY{n}{r\PYZus{}v} \PY{o}{*} \PY{n}{B\PYZus{}max} \PY{o}{*} \PY{n}{R\PYZus{}max} \PY{o}{*} \PY{n}{phi0inv}
    \PY{n}{k\PYZus{}v} \PY{o}{=} \PY{n}{k\PYZus{}val} \PY{o}{+} \PY{n}{dk\PYZus{}v} \PY{o}{/}\PY{l+m+mf}{2.0}\PY{o}{/}\PY{n}{R\PYZus{}max}

    \PY{c+c1}{\PYZsh{} predict a\PYZus{}j, b\PYZus{}j for every row}
    \PY{k}{for} \PY{n}{i} \PY{o+ow}{in} \PY{n+nb}{range}\PY{p}{(}\PY{n}{no\PYZus{}samp}\PY{p}{)}\PY{p}{:}
        \PY{n}{pred\PYZus{}vals\PYZus{}a}\PY{p}{[}\PY{n}{i}\PY{p}{]}\PY{p}{,} \PY{n}{pred\PYZus{}vals\PYZus{}b}\PY{p}{[}\PY{n}{i}\PY{p}{]} \PY{o}{=} \PY{n}{ab\PYZus{}pred\PYZus{}ee}\PY{p}{(}\PY{n}{r\PYZus{}v}\PY{p}{[}\PY{n}{i}\PY{p}{]}\PY{p}{,} \PY{n}{b\PYZus{}v}\PY{p}{[}\PY{n}{i}\PY{p}{]}\PY{p}{,} \PY{n}{dk\PYZus{}v}\PY{p}{[}\PY{n}{i}\PY{p}{]}\PY{p}{,} \PY{n}{mu\PYZus{}v}\PY{p}{[}\PY{n}{i}\PY{p}{]}\PY{p}{,} \PY{n}{dpsi\PYZus{}v}\PY{p}{[}\PY{n}{i}\PY{p}{]}\PY{p}{,} \PY{n}{A\PYZus{}val}\PY{p}{,} \PY{n}{k\PYZus{}val}\PY{p}{)}

    \PY{c+c1}{\PYZsh{} add rows to table}
    \PY{n}{table}\PY{p}{[}\PY{l+s+s2}{\PYZdq{}}\PY{l+s+s2}{ap}\PY{l+s+s2}{\PYZdq{}}\PY{p}{]} \PY{o}{=} \PY{n}{pred\PYZus{}vals\PYZus{}a}
    \PY{n}{table}\PY{p}{[}\PY{l+s+s2}{\PYZdq{}}\PY{l+s+s2}{bp}\PY{l+s+s2}{\PYZdq{}}\PY{p}{]} \PY{o}{=} \PY{n}{pred\PYZus{}vals\PYZus{}b}
\end{Verbatim}
\end{tcolorbox}

    \begin{tcolorbox}[breakable, size=fbox, boxrule=1pt, pad at break*=1mm,colback=cellbackground, colframe=cellborder]
\prompt{In}{incolor}{50}{\boxspacing}
\begin{Verbatim}[commandchars=\\\{\}]
\PY{n}{df} \PY{o}{=} \PY{n}{tbl021}
\PY{n}{predict\PYZus{}table\PYZus{}ab}\PY{p}{(}\PY{n}{df}\PY{p}{)}
\end{Verbatim}
\end{tcolorbox}

    \begin{tcolorbox}[breakable, size=fbox, boxrule=1pt, pad at break*=1mm,colback=cellbackground, colframe=cellborder]
\prompt{In}{incolor}{52}{\boxspacing}
\begin{Verbatim}[commandchars=\\\{\}]
\PY{n}{tbl021}\PY{o}{.}\PY{n}{head}\PY{p}{(}\PY{p}{)}
\end{Verbatim}
\end{tcolorbox}

            \begin{tcolorbox}[breakable, size=fbox, boxrule=.5pt, pad at break*=1mm, opacityfill=0]
\prompt{Out}{outcolor}{52}{\boxspacing}
\begin{Verbatim}[commandchars=\\\{\}]
     R    B    dk            mu  \textbackslash{}
0  0.9  0.5  0.51  1.000000e-10
1  0.9  0.5  0.51  1.584893e-10
2  0.9  0.5  0.51  2.511886e-10
3  0.9  0.5  0.51  3.981072e-10
4  0.9  0.5  0.51  6.309573e-10

                                            dpsi0                  a0  \textbackslash{}
0  1.054613e+10+4.462116e+                    07j  0.501674-0.439413j
1  1.054613e+10+4.462116e+                    07j  0.501674-0.439413j
2  1.054613e+10+4.462116e+                    07j  0.501674-0.439413j
3  1.054613e+10+4.462116e+                    07j  0.501674-0.439413j
4  1.054613e+10+4.462116e+                    07j  0.501674-0.439413j

                   b0   A max 0        I0  \textbackslash{}
0  0.498326+0.439413j  1.331292  0.003348
1  0.498326+0.439413j  1.331292  0.003348
2  0.498326+0.439413j  1.331292  0.003348
3  0.498326+0.439413j  1.331292  0.003348
4  0.498326+0.439413j  1.331292  0.003348

                                             dpsi                   a  \textbackslash{}
0  1.052821e+10+4.460161e+                    07j  0.501673-0.438666j
1  1.051777e+10+4.459022e+                    07j  0.501673-0.438231j
2  1.050128e+10+4.457225e+                    07j  0.501672-0.437544j
3  1.047527e+10+4.454395e+                    07j  0.501671-0.436460j
4  1.043440e+10+4.449953e+                    07j  0.501669-0.434757j

                    b  A max new      I v2  effective mu     I0 p1     I0 p2  \textbackslash{}
0  0.498327+0.438666j   1.330307  0.003346  1.769717e-10  0.003348  0.003348
1  0.498327+0.438231j   1.329733  0.003345  2.802394e-10  0.003348  0.003348
2  0.498328+0.437544j   1.328828  0.003344  4.435448e-10  0.003348  0.003348
3  0.498329+0.436460j   1.327402  0.003342  7.014631e-10  0.003348  0.003348
4  0.498331+0.434757j   1.325164  0.003338  1.107999e-09  0.003348  0.003348

     I pred                  ap                  bp
0  0.003346  0.503126-0.436999j  0.496869+0.440317j
1  0.003345  0.503123-0.436565j  0.496871+0.439881j
2  0.003344  0.503119-0.435879j  0.496875+0.439193j
3  0.003342  0.503113-0.434797j  0.496881+0.438108j
4  0.003338  0.503103-0.433097j  0.496891+0.436402j
\end{Verbatim}
\end{tcolorbox}
        
    \begin{tcolorbox}[breakable, size=fbox, boxrule=1pt, pad at break*=1mm,colback=cellbackground, colframe=cellborder]
\prompt{In}{incolor}{53}{\boxspacing}
\begin{Verbatim}[commandchars=\\\{\}]
\PY{n}{g} \PY{o}{=} \PY{n}{sns}\PY{o}{.}\PY{n}{PairGrid}\PY{p}{(}\PY{n}{tbl021}\PY{p}{[}\PY{p}{[}\PY{l+s+s2}{\PYZdq{}}\PY{l+s+s2}{a}\PY{l+s+s2}{\PYZdq{}}\PY{p}{,} \PY{l+s+s2}{\PYZdq{}}\PY{l+s+s2}{b}\PY{l+s+s2}{\PYZdq{}}\PY{p}{,} \PY{l+s+s2}{\PYZdq{}}\PY{l+s+s2}{ap}\PY{l+s+s2}{\PYZdq{}}\PY{p}{,} \PY{l+s+s2}{\PYZdq{}}\PY{l+s+s2}{bp}\PY{l+s+s2}{\PYZdq{}}\PY{p}{]}\PY{p}{]}\PY{p}{)}

\PY{k}{with} \PY{n}{warnings}\PY{o}{.}\PY{n}{catch\PYZus{}warnings}\PY{p}{(}\PY{n}{action}\PY{o}{=}\PY{l+s+s2}{\PYZdq{}}\PY{l+s+s2}{ignore}\PY{l+s+s2}{\PYZdq{}}\PY{p}{)}\PY{p}{:}
    \PY{n}{g}\PY{o}{.}\PY{n}{map\PYZus{}diag}\PY{p}{(}\PY{n}{sns}\PY{o}{.}\PY{n}{kdeplot}\PY{p}{)}
    \PY{n}{g}\PY{o}{.}\PY{n}{map\PYZus{}upper}\PY{p}{(}\PY{n}{sns}\PY{o}{.}\PY{n}{scatterplot}\PY{p}{)}
    \PY{n}{g}\PY{o}{.}\PY{n}{map\PYZus{}lower}\PY{p}{(}\PY{n}{sns}\PY{o}{.}\PY{n}{scatterplot}\PY{p}{)}
\end{Verbatim}
\end{tcolorbox}

    \begin{center}
    \adjustimage{max size={0.9\linewidth}{0.9\paperheight}}{figs/bvp_fin_en/output_83_0.png}
    \end{center}
    { \hspace*{\fill} \\}
    
    The model is correct. Note that the plots above only show the real
components of \(a_j\) and \(b_j\).

\(a_j+b_j = 1\), hence the inverse proportionality in the plots above
between \(a_j\) and \(b_j\).

    


    % Add a bibliography block to the postdoc
    
    
    
\end{document}
